\documentclass[reqno, 11pt, a4paper]{amsart} 
\usepackage[english]{babel}
\usepackage{amsfonts, amsmath, amsthm, amssymb,amscd,indentfirst}
\usepackage{mathtools}
\usepackage{xcolor}
\colorlet{RED}{red}
%\usepackage[notref,notcite]{showkeys}
%\usepackage{times}
\usepackage{newtxtext,newtxmath}
\usepackage{enumerate}
\usepackage{enumitem}
\usepackage{esint}
\usepackage{stackrel}
\usepackage{float}
%\usepackage{bm}

\usepackage{anysize} %Anysize is a simple package that can be used to set up the margins of a document.
%\usepackage{mathrsfs}
 \usepackage[mathscr]{euscript}
\marginsize{2cm}{2cm}{2cm}{2cm}
%\marginsize{Izque}{Derec}{Arrib}{Abajo}
%\usepackage[shortlabels]{enumitem}

%%%%%%
%%%%%%
%%%%%%%


%\pagecolor{black}
%\color{white}


%%
%%
%%


\usepackage{etoolbox}

\makeatletter

% Líderes de puntos estéticos en el índice
\patchcmd{\@tocline}
  {\hfil}
  {\leaders\hbox to 0.6em{\hss.\hss}\hfill}
  {}{}

% Ajuste del ancho reservado para número de página
\renewcommand{\@pnumwidth}{1.5em}

% Definición de niveles del índice con espaciado e indentación adecuada
\def\l@section{\@tocline{1}{0.6em}{0em}{}{}}         % Sección sin sangría
\def\l@subsection{\@tocline{2}{0.3em}{2em}{}{}}      % Subsección con 2em
\def\l@subsubsection{\@tocline{3}{0.2em}{4.7em}{}{}}   % Subsubsección con 4em

% Corrección para evitar espacio extra innecesario en el índice
\patchcmd{\@tocline}
  {\vskip #1}
  {\vskip #1\relax}
  {}{}

\makeatother


\setcounter{tocdepth}{3}
%%%%%
%%%%%%
%%%%%%%
%%%%%%

%\def\theequation{\thesection.\arabic{equation}}
\newtheorem{theorem}{Theorem}[section]
\newtheorem{proposition}[theorem]{Proposition}
\newtheorem{lemma}[theorem]{Lemma}
\newtheorem{definition}[theorem]{Definition}
\newtheorem{corollary}[theorem]{Corollary}
\newtheorem{example}[theorem]{Example}
\newtheorem{remark}[theorem]{Remark}
\newtheorem{cl}{Claim}
\newtheorem{claim}[theorem]{Claim}
%\newtheorem{thm}{Theorem}[section]
% \newtheorem{cor}[thm]{Corollary}
 %\newtheorem{lem}[thm]{Lemma}
 %\newtheorem{prop}[thm]{Proposition}
 %\theoremstyle{definition}
 %\newtheorem{defn}[thm]{Definition}
 %\theoremstyle{remark}
 %\newtheorem{rem}[thm]{Remark}
% \newtheorem{example}{Example}
 \numberwithin{equation}{section}
 
%\newcommand{\past}{p^\ast _{\alpha , \gamma}}
%\newcommand{\na}{N _{\alpha , \gamma}}
 %\newcommand{\dv}{\textnormal{d}v}
  \newcommand{\dvg}{\textnormal{d}v_g}
 \newcommand{\dx}{\textnormal{d}x}
 \newcommand{\ds}{\textnormal{d}\sigma}
  \newcommand{\dsg}{\textnormal{d}\sigma _g}
 \newcommand{\dt}{\textnormal{d} t}
 \newcommand{\gij}{\textnormal{a}  \textnormal{b}}
 \newcommand{\gcc}{\textnormal{c}  \textnormal{d}}
 %\newcommand{\iwe}{ \mathscr{G}}
 \newcommand{\iwe}{ \tau}
% \newcommand{\iwb}{ \mathscr{G}}
 \newcommand{\ops}{\mathcal{Q}}
 \newcommand{\xv}{2_{\mathtt{V}}}
 \newcommand{\xw}{2_{\mathtt{W}}}
 \newcommand{\qxv}{q_{\mathtt{V}}}
 \newcommand{\qxw}{q_{\mathtt{W}}}
 \newcommand{\ck}{\textnormal K}
 \newcommand{\cty}{\operatorname C _{\nu}}
 \newcommand{\ctya}{\operatorname{Cap}}
 \newcommand{\cne}{C }
 %\newcommand{\cneq}{C_{q ,\mathcal E} }
 \newcommand{\lpm}{\lambda ^{\pm}}
 \newcommand{\fc}{\mathsf c}
\newcommand{\omb}{\Omega}
\newcommand{\wom}{W ^{1, q}(\Omega ; \mathtt{V} _0 , \mathtt{V} _1)}
\newcommand{\womd}{W^{1,2} (\Omega ; \mathtt{V} _0 , \mathtt{V} _1)}
% \renewcommand{\dh}{\textnormal{d}\mathcal{H}^{N-1}}
%  \newcommand{\du}{\textnormal{d}\mu}
 % \newcommand{\dub}{\textnormal{d}\bar{\mu}}
  % \newcommand{\duh}{\textnormal{d}\hat{\mu}}
  %\newcommand{\dun}{\textnormal{d}\mu _n}
% \newcommand{\dvta}{\textnormal{d} \vartheta _\alpha}
 %\newcommand{\dvtao}{\textnormal{d} \vartheta _{\gamma}}
 %\newcommand{\dvtb}{\textnormal{d} \vartheta _\beta}
 %\newcommand{\vta}{\vartheta _\alpha}
  %\newcommand{\vtao}{\vartheta _{\gamma}}
  %\newcommand{\vtb}{\vartheta _\beta}
  \newcommand{\sign}{\operatorname{sign}}
 \newcommand{\spt}{\operatorname{supp}}
% \newcommand{\grf}{F}
\newcommand{\dist}{\operatorname{dist}}
\newcommand{\wghtv}{\mathtt{V}}
\newcommand{\wghtw}{\mathtt{W}}
\newcommand{\dwv}{\textnormal{d} \mathtt{V}}
\newcommand{\dww}{\textnormal{d} \mathtt{W}}
\newcommand{\loc}{\operatorname{loc}}
\newcommand{\kp}{\mathscr{R}}
\newcommand{\rd}{\textnormal{R}}
\newcommand{\dst}{\rho}
\newcommand{\mec}{\delta_{\mathbb{R}^n}}
\newcommand{\diam}{\operatorname{diam}}
\newcommand{\pvr}{p}
\newcommand{\qxp}{q}
\newcommand{\dy}{\textnormal{d}y}
\newcommand{\cen}{\mathtt{C}}
%\newcommand{\soc}{\mathcal{O}}
%\newcommand{\vb}{\mathscr{V}}
\newcommand{\bnorm}{\boldsymbol{\|}}
\newcommand{\bbnorm}{\boldsymbol{\big\|}}


\newcommand{\R}{\mathbb{R}}
\newcommand{\N}{\mathbb{N}}
\newcommand{\op}{\mathrm{op}}
\newcommand{\ind}{\mathbf{1}}

\newcommand{\qxpz}{q_0}
\newcommand{\qxpo}{q_1}
\newcommand{\Bint}{\mathcal{B}_{\qxp,\qxpz}^{m}}
\newcommand{\Bbd}{\mathcal{B}_{\Gamma,\qxp,\qxpo}^{m}}


%\usepackage{xcolor}
%\pagecolor{black}
%\color{white}

\title{{Weighted Sobolev Spaces and an Elliptic Eigenvalue Problem
with an Indefinite Weight on Domains in Noncompact Riemannian Manifolds}}

%\title{{Weighted Sobolev Spaces and an Indefinite Elliptic Eigenvalue Problem on Domains in Noncompact Riemannian Manifolds}}

\author{Juan Pablo Alcon Apaza}


\address{Juan Pablo Alcon Apaza, Departamento de Matem\'atica, Universidade Federal de Minas Gerais, 31270-901, Belo Horizonte - MG,  Brazil}

\email{juanpabloalconapaza@gmail.com}







\begin{document}
\maketitle
\begin{abstract} The aim of this work is to establish sufficient conditions ensuring
the continuity and compactness of the weighted Sobolev embedding and trace operators $$ W^{1,\qxp}(\Omega; \wghtv_0, \wghtv_1) \rightarrow L^{\qxp_0}(\Omega; \wghtv_2) \quad \text { and } \quad W^{1,\qxp}(\Omega; \wghtv_0, \wghtv_1) \rightarrow L^{\qxp_1}(\partial \Omega; \wghtw).$$ As an application, we study the Neumann eigenvalue problem with an indefinite weight \begin{equation*}\left\{\begin{aligned}-\operatorname{div} ( \wghtv _1 \nabla u )+ \wghtv _0 u   & = \lambda \wghtv _2 \iwe u  & & \text { in } \Omega,\\\wghtv _1 \frac{\partial u}{\partial \nu} &= 0 & & \text { on } \partial \Omega .\end{aligned}\right.\end{equation*} where $\Omega$ is an open subset of a noncompact Riemannian manifold, $\lambda$ is a real number, $\iwe$ is a sign-changing function, and $\wghtv _0$, $\wghtv _1$, \(\wghtv_2\), and \(\wghtw\) are weight functions satisfying suitable conditions.  We prove that this problem has infinitely many positive and negative eigenvalues. We also establish boundedness of its weak eigenfunctions and, for $u\in W^{1,2}_0(\overline{\Omega};\wghtv_0,\wghtv_1)$, the decay property  $$\lim_{m\rightarrow\infty}\operatorname*{ess\,sup}_{\Omega\setminus\overline{D_m}}|u|=0,$$ where $(D_m)_{m\in\mathbb N}$ is an exhaustion of the manifold by bounded open sets.\end{abstract}




\let\thefootnote\relax\footnote{2020 \textit{Mathematics Subject Classification}.  {46E35; 35J25; 35P05}}
\let\thefootnote\relax\footnote{\textit{Keywords and phrases}. Weighted Sobolev spaces, compact embeddings, trace operators, indefinite eigenvalue problems, noncompact Riemannian manifolds}

%35J25 Boundary value problems for second-order elliptic equations
%35P05 General topics in linear spectral theory for PDEs
%46E35 Sobolev spaces and other spaces of "smooth" functions, embedding theorems, trace theorems

\markright{{WEIGHTED SOBOLEV SPACES AND AN INDEFINITE ELLIPTIC EIGENVALUE PROBLEM}}


\tableofcontents
\section{Introduction}



Weighted Sobolev spaces are widely used as solution spaces for degenerate elliptic equations. Kufner \cite{kugner1985weightedsobo}, Triebel \cite{triebel1995interpolati}, and Schmeisser \& Triebel \cite{schmeisser1987topics}, among others, have contributed to their study on unbounded domains. See, for example, \cite{choquet2008einsteinequations, chua1992extensiontheorems, gruka1991weighteiii} for the Euclidean case, \cite{hamann2013singularmanifolds, pacini2013conicalsing} for noncompact Riemannian domains, and \cite{zbMATH06237307, post2023heatkernel} for unweighted Sobolev spaces on noncompact Riemannian manifolds.

In the study of partial differential equations on a domain \(\Sigma\subset\mathbb R^n\), it is often useful to know that the embedding of \(W^{m,\qxp}(\Sigma)\) into \(L^{\qxp_0}(\Sigma)\) is compact. For nonlinear variational equations, this property is often used to prove that the associated energy functional satisfies the Palais--Smale condition. When \(\Sigma\) is unbounded, compactness generally fails. Sufficient conditions involving weight functions can be found, for example, in \cite{adams1971compactembedding, gruka1991weighteiii, kufner1990hardytype}.

{ 
The elliptic application considered in this paper is the homogeneous indefinite eigenvalue problem
}
\begin{equation}
\left\{
\begin{aligned}
-\operatorname{div} ( \wghtv _1 \nabla u )+ \wghtv _0 u   & = \lambda \wghtv _2 \iwe u & & \text { in } \Omega,\\
\wghtv _1 \frac{\partial u}{\partial \nu} &= 0 & & \text { on } \partial \Omega .
\end{aligned}
\right. \label{61}
\end{equation}
Here, $\Omega$ is an open subset of a Riemannian manifold $M$ of dimension $n\geq2$, $\nu$ is the outward unit normal vector to $\partial\Omega$, $\lambda$ is a real number, $\iwe:M\rightarrow\mathbb R$ is a sign-changing function, and $\wghtv_0$, $\wghtv_1$, and $\wghtv_2$ are weight functions satisfying suitable conditions.

{ 
We also determine conditions under which weak solutions of \eqref{61} are bounded and satisfy
\begin{equation} \label{87}
\lim_{m\rightarrow\infty}\stackbin[\Omega\setminus\overline{D_m}]{}{\operatorname{ess}\sup}\,|u|=0.
\end{equation}
}
Here, $(D_m)_{m\in\mathbb N}$ is an increasing sequence of bounded open subsets of $M$ such that
$$
M=\bigcup_m D_m, \qquad D_m\Subset D_{m+1},
$$
and $D_m\cap\Omega$ has Lipschitz boundary for every $m\in\mathbb N$. We also define
$$
D^m:=M\setminus\overline{D_m}.
$$







\subsection{Context and Related Work}

{On a closed Riemannian manifold, the Laplace--Beltrami operator has discrete spectrum. For an open subset of $\mathbb R^n$, discreteness depends on the chosen realization of the Laplacian and on the boundary condition. In particular, the Neumann Laplacian has discrete spectrum whenever the embedding $W^{1,2}(\mathcal M)\rightarrow L^2(\mathcal M)$ is compact.}
In general, the spectrum of $\Delta_{\mathcal M}$ need not be discrete. Several nonstandard geometric settings have therefore been studied. Conditions for discreteness of the spectrum of the Laplacian on noncompact complete Riemannian manifolds with specific geometric structures appear in \cite{baider1979discretespectra, brooks1984fiitevolume, bruning1989discretespectrum, donnelly1979purepointspectrum, escobar1986spectrum, kleine1988discretenessconditions}.

Cianchi \& Maz'ya \cite{cianchi2011spectrumnoncomriem} prove that, for a noncompact Riemannian manifold $\mathcal M^n$ with $n\geq2$ and $\mathcal H^n(\mathcal M)<\infty$, the embedding $W^{1,2}(\mathcal M)\rightarrow L^2(\mathcal M)$ is compact if and only if
$$
\lim_{s\rightarrow0}\frac{s}{\mu_{\mathcal M}(s)}=0,
$$
which is equivalent to discreteness of the spectrum of $\Delta_{\mathcal M}$. Here, the isocapacitary function $\mu_{\mathcal M}:[0,\mathcal H^n(\mathcal M)/2]\rightarrow[0,\infty]$ is defined by
\begin{align*}
\mu_{\mathcal M}(s):=\inf\Big\{C(E,G)\:|\: E \text { and } G & \text { are measurable subsets of } \mathcal M,\\
& E\subset G\subset\mathcal M,\quad s\leq\mathcal H^n(E),\quad \mathcal H^n(G)\leq\frac12\mathcal H^n(\mathcal M)\Big\},
\end{align*}
where
\begin{align*}
  C(E, G):=\inf \Bigg\{  \int_{\mathcal M}|\nabla u|^2 \, \dx \:|\: u \in W^{1,2}(M), & u \geq 1  \text  { in } E \text  {  and }  u \leq 0\\
  & \text { in }  \mathcal M \setminus G \text  { (up to a set of standard capacity zero) }\Bigg\}.
  \end{align*}


Furthermore, Cianchi \& Maz'ya \cite{cianchi2013boundseigenfunctions} prove that, if $\mathcal H^n(\mathcal M)<\infty$ and
$$
\int_0\frac{\textnormal{d}s}{\mu_{\mathcal M}(s)}<\infty,
$$
then, for every eigenvalue $\gamma$ of $\Delta_{\mathcal M}$, there exists a constant $C=C(\mu_{\mathcal M},\gamma)$ such that
$$
\|u\|_{L^\infty(\mathcal M)}\leq C\|u\|_{L^2(\mathcal M)}
$$
for every eigenfunction $u$ associated with $\gamma$.



Skrzypczak \& Tintarev \cite{zbMATH07331136} study compact subsets of Sobolev spaces on complete, noncompact, connected Riemannian manifolds. They define orbital discretizations and functions that are quasisymmetric relative to an orbital discretization, and use the spotlight lemma to prove that a bounded set
\[
K\subset H^{1,p}(M)\cap S_{\Gamma,i,\lambda}(M)
\]
is relatively compact in \(L^q(M)\). For a compact, connected group \(G\) of isometries, they prove that if \(G\) is coercive, then \(H_G^{1,p}(M)\) is compactly embedded into \(L^q(M)\). Conversely, compactness implies that \(G\) is coercive provided that the injectivity radius of \(M\) is positive.

%Esto reproduce el contenido de los teoremas 3.5 y 4.3, incluyendo la condición necesaria sobre el radio de inyectividad. 


%Skrzypczak \& Tintarev \cite{zbMATH07331136} study compactness phenomena for Sobolev embeddings on complete, connected, noncompact Riemannian manifolds. They introduce orbital discretizations and a corresponding notion of quasisymmetry, and use the spotlight lemma to prove that bounded sets of quasisymmetric Sobolev functions are relatively compact in subcritical Lebesgue spaces.  For compact, connected groups of isometries, they show that coercivity of the group action implies compactness of the invariant Sobolev embedding. %Conversely, compactness implies coercivity when the manifold has positive injectivity radius; thus, under this additional assumption, compactness is characterized by coercivity of the action. 





{ 
Let $\mathcal M_0$ be a complete Riemannian manifold, and let $\mathsf m\geq0$ be a locally integrable potential. Ouhabaz \cite{ouhabaz2001spectral} studies the positivity of the $L^2$ spectral lower bound of the Schrödinger operator $-\Delta_{\mathcal M_0}+\mathsf m$.
}
He proves that, if $\mathcal M_0$ satisfies local $L^2$-Poincaré inequalities, namely,
$$
\int_{B(\pvr,r)}\left|u-\overline{u_{B(\pvr,r)}}\right|^2\,\dvg
\leq C(R)r^2\int_{B(\pvr,r)}|\nabla u|^2\,\dvg
$$
for every $u\in C^\infty(B(\pvr,r))$, $\pvr\in\mathcal M_0$, and $0<r\leq R$, and if $\mathcal M_0$ satisfies a local doubling property, then
$$
\inf\sigma(-\Delta_{\mathcal M_0}+\mathsf m)>0,
$$
provided that
$$
\inf_{\pvr\in\mathcal M_0}\frac{1}{|B(\pvr,r)|}\int_{B(\pvr,r)}\mathsf m\,\dvg>0
$$
for some $r>0$. Ouhabaz also shows that this mean condition is necessary under additional geometric assumptions on $\mathcal M_0$.

Berestycki \& Rossi \cite{zbMATH06441399} study three definitions of the generalized principal eigenvalue for linear second-order elliptic operators on unbounded domains. They show that several classical properties of the Dirichlet principal eigenvalue may fail in this setting. They also relate these generalized eigenvalues to the maximum principle and to the existence of positive eigenfunctions satisfying Dirichlet boundary conditions. Since the Dirichlet resolvent is not compact in general, they do not rely on the Krein--Rutman theorem. Instead, they use positive sub- and supersolutions, exhaustion by bounded domains, comparison principles, Harnack inequalities, and local elliptic estimates.



If $\Sigma_0$ is a bounded domain in $\mathbb R^n$ with smooth boundary, the eigenvalue problem
$$
\left\{
\begin{aligned}
-\Delta u &=\lambda u & & \text { in } \Sigma_0,\\
u &=0 & & \text { on } \partial\Sigma_0,
\end{aligned}
\right.
$$
has an infinite sequence of positive eigenvalues
$$
0<\lambda_1<\lambda_2\leq\cdots\leq\lambda_k\leq\cdots,\qquad \lambda_k\rightarrow\infty \quad\text { as } k\rightarrow\infty,
$$
each of finite multiplicity. The same properties hold for
$$
\left\{
\begin{aligned}
-\Delta u+a_0(x)u &=\lambda m(x)u & & \text { in } \Sigma_0,\\
u &=0 & & \text { on } \partial\Sigma_0,
\end{aligned}
\right.
$$
when $a_0$ and $m$ are positive and sufficiently regular on $\overline{\Sigma_0}$. For a detailed study of principal eigenvalues for second-order differential operators that are not necessarily in divergence form, see Fleckinger, Hernández, \& Thélin \cite{fleckinger2004existence}.

The existence of principal eigenvalues with positive eigenfunctions is important in nonlinear problems for which positive solutions are sought. This occurs, for example, in reaction--diffusion systems arising in population dynamics, chemical reactions, and combustion; see \cite{smoller2012shock}.

The classical reference for this theory is the book by Courant \& Hilbert \cite{courant1962methods}. The main tool is the spectral theory of compact self-adjoint operators on Hilbert spaces. { The variational characterization yields monotonicity with respect to the coefficients and, under domain inclusion, with respect to the domain. Continuity with respect to coefficients or domains requires an appropriate notion of convergence.}

{ 
The theory also extends to unbounded coefficients and indefinite weights. Assume that $a_0>0$ in $\Sigma_0$, $a_0,m\in L^r(\Sigma_0)$ for some $r>n/2$, and that the positive part of $m$ is nontrivial. Then there exists a unique positive principal eigenvalue $\lambda_1^+>0$ with a positive eigenfunction. If the negative part of $m$ is nontrivial, there is also a unique negative principal eigenvalue $\lambda_1^-<0$ with a positive eigenfunction; see \cite{figuereido1982semilinearelliptic, brown1980existence, weinberger1974variational}.
}




%%%%%%%%%%%%%%%%%%%%%%%%%%%%%%%%%%%%%%%%%%%%%%%%%%%%%%%%%%%%%%%%%%%%%%%%%%%%%%%%%%%%%%%%%%%%%%%%%%%%%%%%%%%%%%%%%%%%%%%%%%%%%%%%%%%%%%%%%%%%%%%%%%%%%%%%%%%%%%%%%%%%%%%%%%%%%%%%%%%%%%%%%%%%%%%%%%%%%%%%%%%%%%%%%%%%%%%%%%%%%%%%%%%%%%%%%%%%%%%%%%%%%%%%%%%%%%%%%%%%%%%%%%%%%%%%%%%%%%%%%%%%%%%%%%%%%%%%%%%%%%%%%%%%%%%%%%%%%%%%%%%%%%%%%%%%%%%%%%%%%%%%%%%%%%%%%%%%%%%%%%%%%%%%%%%%%%%%%%%%%%%%%%%%%%%%%%%%%%%%%%%%%%%%%%%%%%%%%%%%%%%%%%%%%%%%%%%%%%%%%%%%%%%%%%%%%%%%%%%%

%Regarding elliptic problems involving \(L^1\) or measure data, the main mathematical difficulty consists in the fact that the classical variational formulation is not possible. Existence results have been obtained using a non-variational framework for these problems. The first approach is due to Stampacchia \cite{stampacchia1965dirichlet}, who obtained solutions with a duality method. However, a limitation of this approach is that it applies only to linear equations. In \cite{boccardogallouet1989measuredata}, Boccardo and Gallouët have developed a method based on the approximation of the data by smoother functions. Applications of these two methods in the study of several types of elliptic problems with Dirichlet boundary conditions and \(L^1\) or measure data may be found in \cite{bocearadulescu1997elliptiques, murat1993renormalized, orsina1993solvabilityl1data, prignet1995elliptic, radulescu1996neutrons}.

\subsection{Main Results} In this work, we extend the results of Pflüger \cite{pfluger1998compactraces} to noncompact Riemannian manifolds. Under suitable conditions on the weights and on the geometry of \(\Omega\), we establish continuous and compact embeddings and traces for \(W^{1,\qxp}(\Omega;\wghtv_0,\wghtv_1)\); see conditions \ref{5}--\ref{3} below. For other geometric and weighted settings, see, for example, \cite{hamann2013singularmanifolds}, where Sobolev--Slobodeckii and Bessel potential spaces on singular manifolds are studied.

We also obtain { two sequences of eigenvalues} for problem \eqref{61}, following the method of Allegretto \cite{allegreto1992eigenvalueindefiniteweight}. The main difficulty is to prove the decay property \eqref{87} on the unbounded domain \(\Omega\). For Fredholm properties of elliptic operators on noncompact manifolds, see Lockhart \& McOwen \cite{lockartowen1985ellipticoperators} and Lockhart \cite{lockhart1980fredholm}.

%Following the approaches of Bocea \& Redulescu \cite{bocea1999eigenvaluproblemweight} and Orsina \cite{orsina1993solvabilityl1data}, we will then seek weak solutions of problem \eqref{66} that admit $L^1$ data, after having studied problem \eqref{61}.


 
 



%%%%%%%%%%%%%%%%%%%%%%%%%%%%%%%%%%%%%%%%%%%%%%%%%%%%%%%%%%%%%%%%%%%%%%%%%%%%%%%%%%%%%%%%%%%%%%%%%%%%%%%%%%%%%%%%%%%%%%%%%%%%%%%%%%%%%%%%%%%%%%%%%%%%%%%%%%%%%%%%%%%%%%%%%%%%%%%%%%%%%%%%%%%%%%%%%%%%%%%%%%%%%%%%%%%%%%%%%%%%%%%%%%%%%%
  

In the sequel, we write \(\Gamma=\partial\Omega\). We assume that there exists {a locally finite family of open sets \(U_{k,i},\hat U_{k,i}\subset M\), \((k,i)\in\{0,1\}\times\mathbb N\), covering a neighborhood of \(\overline\Omega\)}, with the following properties:

\begin{enumerate}[label=($U_{\arabic*}$)]
\item \label{5} The family satisfies
\[
\overline{\Omega}
\subset
\left(\bigcup_i U_{0,i}\right)\cup\left(\bigcup_j U_{1,j}\right),
\quad
\bigcup_i\hat U_{0,i}\subset\Omega,
\quad
\Gamma\subset\bigcup_j U_{1,j}, \quad  M= \bigcup _{(k,i)\in \{0,1\} \times \mathbb{N}} \hat U _{k,i},
\]
and \(U_{k,i}\subset\hat U_{k,i}\) for every \((k,i)\in\{0,1\}\times\mathbb N\). %{\color{red}Thus, \(\bigcup_{k,i}\hat U_{k,i}\) is an open neighborhood of \(\overline\Omega\).}

\item \label{2} For every \((k,i)\in\{0,1\}\times\mathbb N\), there exist \(0<r_{k,i}<\hat r_{k,i}\) and a chart
\[
\psi_{k,i}:B(0,\hat r_{k,i})\subset\mathbb R^n\longrightarrow\hat U_{k,i}\subset M
\]
such that:
\begin{enumerate}[label=($\alph*$)]
   \item The restriction \(\psi_{k,i}:B(0,r_{k,i})\rightarrow U_{k,i}\) is a diffeomorphism.

   \item For every \(j\in\mathbb N\),
   \[
   { 
   \psi_{1,j}\left(B(0,\hat r_{1,j})\cap\{x_n>0\}\right)
   =\hat U_{1,j}\cap\Omega,
   \qquad
   \psi_{1,j}\left(B(0,\hat r_{1,j})\cap\{x_n=0\}\right)
   =\hat U_{1,j}\cap\Gamma.
   }
   \]

   \item For every \(j\in\mathbb N\), the map
   \[
   \psi_{\Gamma,1,j}:B(0,\hat r_{1,j})\cap\{x_n=0\}\longrightarrow\hat U_{1,j}\cap\Gamma,
   \]
   defined by
   \[
   \psi_{\Gamma,1,j}(x_1,\ldots,x_{n-1})
   =\psi_{1,j}(x_1,\ldots,x_{n-1},0),
   \]
   is a chart for \(\Gamma\).
\end{enumerate}
% \item \label{32} There exist $\rd _1 >0$ and $\ck _2 >0$ such that $r_{k,i}< \rd _1 $ and , for every $(k,i)\in \{0,1\} \times \mathbb{N}$.

\item \label{1} There exists a constant \(\rd_1>0\) such that
\[
\sum_{k,i}\chi_{\hat U_{k,i}}\leq\rd_1
\quad\text{in }M.
\]

\item \label{3} There exists a constant \(\rd_2>0\) such that, for every \(j\in\mathbb N\),
\[
\left(
\sup_{\substack{x\in B(0,\hat r_{1,j})\\ |z|_{\delta_{\mathbb R^n}}=1}}
\left|d(\psi_{1,j})_x z\right|_g
\right)
\left(
\sup_{\substack{(\pvr,v)\in T\hat U_{1,j}\\ |v|_g=1}}
\left|d(\psi_{1,j}^{-1})_{\pvr}v\right|_{\delta_{\mathbb R^n}}
\right)
\leq\rd_2
\]
and
\[
\left(\sup_{\hat U_{1,j}}\det[g_{\gij}]\right)
\left(\sup_{\hat U_{1,j}}\det[g^{\gij}]\right)
\leq\rd_2^2.
\]
Here, \(\delta_{\mathbb R^n}\) is the Euclidean metric,
\([g^{\gij}]:=[g_{\gij}]^{-1}\), and
\(g_{\gij}=g(d\psi_{1,j}e_{\textnormal a},d\psi_{1,j}e_{\textnormal b})\).
\end{enumerate}



\begin{remark} $ $
\begin{enumerate}[label=(\roman*)] 
%\item The number $\rd _1 $ is related to the Besicovitch Covering Theorem, see \cite{benedetto2002real, heinonen2001analysismetric}.


\item Examples satisfying conditions \ref{5}--\ref{3} can be constructed in the hyperbolic ball \(\mathbb B^n\). In Euclidean space, examples include cylinders and the half-space.

\item Conditions \ref{5}--\ref{3} extend to Riemannian manifolds the conditions U1--U4 in \cite{pfluger1998compactraces}; see also \cite{zbMATH01109876}. These conditions control the geometry of \(\partial\Omega\) at infinity.

\item \label{109} In conditions \ref{5}--\ref{3}, the balls \(B(0,r_{1,j})\) and \(B(0,\hat r_{1,j})\) may be replaced by the cubes
\[
Q(0,r_{1,j})
=\{x\in\mathbb R^n\mid\|x\|_{\max}<r_{1,j}\}
\]
and
\[
Q(0,\hat r_{1,j})
=\{x\in\mathbb R^n\mid\|x\|_{\max}<\hat r_{1,j}\},
\]
without changing the results.

\item If \((M,\Omega)\) has {\it bounded geometry}; see \cite[Definition 4.1]{zbMATH06237307} and \cite[Proposition 3.2]{zbMATH01606240}, then conditions \ref{5}--\ref{3} hold for a suitable family of {\it geodesic normal coordinates} and {\it Fermi coordinates}.





\end{enumerate}
\end{remark}


%%%%
%%%%
%%%%
%%%%



%%%%%%%%%%%%%%%%%%%%%%%%%%%%%%%%%%%%%%%%%%%%%%%%%%%%%%%%%%%%%%%%%%%%%%%%%%%%%%%%%%%%%%%%%%%%%



Let \(\wghtv_i\), \(i=0,\ldots,3\), and \(\wghtw\) be weights on \(M\); that is, they are locally integrable and positive almost everywhere. We also assume that \(\wghtw\) is continuous and that all the weights are bounded above and below by positive constants on every compact subset of \(M\).


%%
%%
%%
%%
%%


We assume that there exist positive, continuous functions $b_i:M\rightarrow\mathbb R$, $i=1,2,3$, and constants $\ck_0,\ck_{\qxp}>0$ such that the following conditions hold.

If $(k,i)\in\{0,1\}\times\mathbb N$ and $j\in\mathbb N$, then
\begin{enumerate}[label=($\mathtt{W}_{\arabic*}$)]

\setcounter{enumi}{-1}

 

\item \label{138}
For every $x\in B(0,\hat r_{1,j})\cap\{x_n<0\}$ and every $\ell\in\{0,1\}$,
$$
\wghtv_\ell(\psi_{1,j}(x))
\leq
\ck_0\wghtv_\ell(\psi_{1,j}(x_1,\ldots,x_{n-1},-x_n)).
$$

\item \label{45}
For almost every $\pvr\in\hat U_{k,i}$,
$$
\kp_{k,i}^{\qxp}(\pvr)\wghtv_1(\pvr)
\leq
\ck_{\qxp}\wghtv_0(\pvr).
$$


\item \label{139} For almost every $\pvr\in\hat U_{k,i}$,
$$
\hat r_{k,i}^{-\qxp}\bnorm d\psi_{k,i}\bnorm^{-\qxp}\wghtv_1(\pvr)
\leq
\ck_{\qxp}\wghtv_0(\pvr).
$$


\item \label{46}
For almost every $\pvr\in\hat U_{k,i}$,
$$
\wghtv_2(\pvr)
\leq
b_2(\psi_{k,i}(0))
\quad\text{and}\quad
\bnorm d\psi_{k,i}\bnorm^{\qxp} b_1(\psi_{k,i}(0))
\leq
\wghtv_1(\pvr).
$$

\item \label{16}
For almost every $\pvr\in\hat U_{1,j}$,
$$
\wghtw(\pvr)
\leq
b_3(\psi_{1,j}(0))
\quad\text{and}\quad
{ \bnorm d\psi_{1,j}\bnorm^{\qxp}} b_1(\psi_{1,j}(0))
\leq
\wghtv_1(\pvr).
$$
\end{enumerate}




%%%
%%
%%
%%
%%%
%%

Here,
\begin{align*}
\kp_{k,i}
&\displaystyle:=
\sum_{(\textnormal a,\textnormal b)\in\{0,1\}\times\mathbb N}
\bnorm d\psi_{\textnormal a,\textnormal b}^{-1}\bnorm
(\hat r_{\textnormal a,\textnormal b}-r_{\textnormal a,\textnormal b})^{-1}
\chi_{\hat U_{\textnormal a,\textnormal b}\cap\hat U_{k,i}},\\[5pt]
\bnorm d\psi_{k,i}\bnorm
&\displaystyle:=
\sup\left\{
|d(\psi_{k,i})_x z|_g
\mid
x\in B(0,\hat r_{k,i}),\ |z|_{\delta_{\mathbb R^n}}=1
\right\},\\[5pt]
\bnorm d\psi_{k,i}^{-1}\bnorm
&\displaystyle:=
\sup\left\{
|d(\psi_{k,i}^{-1})_{\pvr}v|_{\delta_{\mathbb R^n}}
\mid
(\pvr,v)\in T\hat U_{k,i},\ |v|_g=1
\right\}.
\end{align*}

Next, set
\begin{equation*}
\mathcal B_{\qxp,\qxp_0}^m
:=
\sup_{\substack{(k,i)\in\{0,1\}\times\mathbb N\\ \psi_{k,i}(0)\in D^m}}
\frac{b_2^{1/\qxp_0}(\psi_{k,i}(0))}{b_1^{1/\qxp}(\psi_{k,i}(0))}
\bnorm G_{k,i}\bnorm^{1/\qxp_0}
\bnorm G_{k,i}^{-1}\bnorm^{1/\qxp}
\hat r_{k,i}^{\frac n{\qxp_0}-\frac n{\qxp}+1},
\end{equation*}
where
\[
\bnorm G_{k,i}\bnorm
:=\sup_{\hat U_{k,i}}\sqrt{\det[g_{\gij}]}
\quad\text{and}\quad
\bnorm G_{k,i}^{-1}\bnorm
:=\sup_{\hat U_{k,i}}\sqrt{\det[g^{\gij}]}.
\]

Furthermore, define
\begin{equation*}
\mathcal B_{\Gamma,\qxp,\qxp_1}^m
:=
\sup_{\substack{j\in\mathbb N\\ \psi_{1,j}(0)\in\Gamma^m}}
\frac{b_3^{1/\qxp_1}(\psi_{1,j}(0))}{b_1^{1/\qxp}(\psi_{1,j}(0))}
\bnorm G_{\Gamma,1,j}\bnorm^{1/\qxp_1}
\bnorm G_{1,j}^{-1}\bnorm^{1/\qxp}
\hat r_{1,j}^{\frac{n-1}{\qxp_1}-\frac n{\qxp}+1},
\end{equation*}
where
\[
{ 
\bnorm G_{\Gamma,1,j}\bnorm
:=
\sup_{\hat U_{1,j}\cap\Gamma}\sqrt{\det[g_{\Gamma,\gij}]}
}
\quad\text{and}\quad
g_{\Gamma,\gij}
:=g(d\psi_{\Gamma,1,j}e_{\textnormal a},d\psi_{\Gamma,1,j}e_{\textnormal b}).
\]




%%
%%
%%
%%
%%
%%




%%
%%
%%
%%
%%
%%
%%
%%



\begin{remark} $ $
\begin{enumerate}[label=(\roman*)] 

\item The quantities \(\mathcal B_{q,q_0}^m\) and \(\mathcal B_{\Gamma,q,q_1}^m\)  extend to Riemannian manifolds the quantities \(\mathcal B_{n,k}\) in \cite{pfluger1998compactraces} and \(\mathscr A_n\) in \cite{gruka1991weighteiii}, which are defined in the Euclidean setting.

%\item Conditions \ref{138} and \ref{45} are used to construct the extension operator in Proposition \ref{22}. They may be omitted if \(\Omega\) already admits a bounded extension operator and the second inequality in \ref{45} holds.


\item If \((M,\Omega)\) has {\it bounded geometry}; see \cite[Definition 4.1]{zbMATH06237307} and \cite[Proposition 3.2]{zbMATH01606240}, then one may choose geodesic normal and Fermi coordinates such that { there exists \(c\geq1\) for which each of
\[
\kp_{k,i},\quad
\bnorm G_{k,i}\bnorm,\quad
\bnorm G_{k,i}^{-1}\bnorm,\quad
\bnorm G_{\Gamma,1,j}\bnorm,\quad
\bnorm d\psi_{k,i}\bnorm,\quad
\bnorm d\psi_{k,i}^{-1}\bnorm,\quad
\hat r_{k,i},\quad
\hat r_{k,i}-r_{k,i}
\]
is bounded above by \(c\) and below by \(c^{-1}\)}.
\end{enumerate}
\end{remark}





%%
%%
%%
%%
%%
%%
%%


\begin{example}\label{ex:shrinking-components-eta}
Let $\eta>0$. Let $\ell_0 > 4^{1/\eta}$. For every
$\ell\in\N$ with $\ell\geq\ell_0$, set
\[
a_\ell:=(\ell,0),
\qquad
R_\ell:=\left( \frac1\ell \right)^\eta,
\qquad
\rho_\ell:=\sqrt{R_\ell}.
\]
Consider the open set
\[
{ 
\Omega
:=
\bigcup_{\substack{\ell\in\N\\ \ell\geq\ell_0}}B(a_\ell,R_\ell)
\subset\R^2.
}
\]
Endow \(\mathbb R^2\) with the Euclidean metric \(g=\delta_{\R^2}\). { For this example, take \(D_m=B(0,m)\).} % Increasing $\ell_0$ if necessary, we may assume that the sets $B(a_\ell,2R_\ell)$, $\ell\in\N$, $\ell\geq\ell_0$, are pairwise disjoint.

Define
\[
r_{0,\ell}:=\left(1-\frac{R_\ell}{4}\right)\rho_\ell,
\qquad
\hat r_{0,\ell}:=\left(1-\frac{R_\ell}{8}\right)\rho_\ell,
\]
and
\[
r_{1,\ell}:=\frac{3\rho_\ell}{4},
\qquad
\hat r_{1,\ell}:=\frac{7\rho_\ell}{8}.
\]

The interior chart is
\[
\psi_{0,\ell}(x,y):=a_\ell+\rho_\ell(x,y),
\qquad
(x,y)\in B(0,\hat r_{0,\ell}).
\]

For the boundary charts, put
\[
S_\ell(x):=\sqrt{R_\ell^2-\rho_\ell^2x^2}.
\]
Define
\[
\psi_{1,\ell}^{\pm }(x,y)
:=
a_\ell+
\left(1-\frac{y}{\rho_\ell}\right)
\left(
\rho_\ell x,
\pm S_\ell(x)
\right),
\qquad
(x,y)\in B(0,\hat r_{1,\ell}),
\]
and
\[
\varphi_{1,\ell}^{\pm}(x,y)
:=
a_\ell+
\left(1-\frac{y}{\rho_\ell}\right)
\left(
\pm S_\ell(x),
\rho_\ell x
\right),
\qquad
(x,y)\in B(0,\hat r_{1,\ell}).
\]



After relabeling the charts
\[
\psi_{0,\ell},
\qquad
\psi_{1,\ell}^{+},
\quad
\psi_{1,\ell}^{-},
\quad
\varphi_{1,\ell}^{+},
\quad
\varphi_{1,\ell}^{-},
\qquad
\ell\in\N,\ \ell\geq\ell_0,
\]
as a family $\psi_{k,i}$, the corresponding covering satisfies conditions
\ref{5}--\ref{3}. Moreover, there exists a constant $C>0$, independent of
$\ell$, such that, for every $\ell\in\N$ with $\ell\geq\ell_0$,
\begin{gather*}
\bnorm d\psi_{0,\ell}\bnorm, \
\bnorm d\psi_{1,\ell}^{\pm}\bnorm,  \
\bnorm d\varphi_{1,\ell}^{\pm}\bnorm
\leq
C\rho_\ell,
\label{eq:dpsi-shrinking-components-eta}
\\[3pt]
\bnorm d(\psi_{0,\ell})^{-1}\bnorm, \
\bnorm d(\psi_{1,\ell}^{\pm})^{-1}\bnorm, \
\bnorm d(\varphi_{1,\ell}^{\pm})^{-1}\bnorm
\leq
\frac{C}{\rho_\ell},
\label{eq:inverse-dpsi-shrinking-components-eta}
\\[3pt]
\bnorm G_{\psi_{0,\ell}}\bnorm,
\bnorm G_{\psi_{1,\ell}^{\pm}}\bnorm, \
\bnorm G_{\varphi_{1,\ell}^{\pm}}\bnorm
\leq
CR_\ell, 
\qquad
\bnorm G_{\psi_{0,\ell}}^{-1}\bnorm, \
\bnorm G_{\psi_{1,\ell}^{\pm}}^{-1}\bnorm, \
\bnorm G_{\varphi_{1,\ell}^{\pm}}^{-1}\bnorm
\leq
\frac{C}{R_\ell},
\label{eq:G-shrinking-components-eta}
\\[3pt]
\bnorm G_{\Gamma,\psi_{1,\ell}^{\pm}}\bnorm, \
\bnorm G_{\Gamma,\varphi_{1,\ell}^{\pm}}\bnorm
\leq
C\rho_\ell, 
\label{eq:Gamma-G-shrinking-components-eta}\\[3pt]
\kp_{\psi_{0,\ell}}, \
\kp_{\psi_{1,\ell}^{\pm}}, \
\kp_{\varphi_{1,\ell}^{\pm}}
\leq
\frac{C}{R_\ell^2} .
\label{eq:kappa-shrinking-components-eta}
\end{gather*}

Define the weights
\[
\wghtv_0(p)
:=
(1+|p|)^{2\eta\qxp+\alpha_0},
\qquad
\alpha_0\geq0,
\]
\[
\wghtv_1(p):=1,
\]
\[
\wghtv_2(p)
:=
(1+|p|)^{-\alpha_2},
\qquad
\alpha_2\geq0,
\]
and
\[
\wghtw(p)
:=
(1+|p|)^{-\beta},
\qquad
\beta\geq0.
\]

There exist positive continuous functions $b_1$, $b_2$, and $b_3$ such that
\[
b_1(x,y)\approx1,
\qquad
b_2(x,y)\approx(1+|(x,y)|)^{-\alpha_2},
\qquad
b_3(x,y)\approx(1+|(x,y)|)^{-\beta}.
\]

Then the weights satisfy conditions \ref{138}--\ref{16}.

Assume that \(\qxp_0\geq\qxp\geq1\) and
\[
-\frac{\alpha_2}{\qxp_0}
-\frac{2\eta}{\qxp_0}
+\frac{2\eta}{\qxp}
-\frac{\eta}{2}
\leq0.
\]
Then there exist constants \(C_1>0\) and \(m_0\in\N\) such that, for every \(m\in\N\) with \(m\geq m_0\),
\[
\mathcal{B}_{\qxp,\qxp_0}^m
\leq
C_1
m^{
-\frac{\alpha_2}{\qxp_0}
-\frac{2\eta}{\qxp_0}
+\frac{2\eta}{\qxp}
-\frac{\eta}{2}
}.
\]


Similarly, assume that $\qxp_1\geq1$ and
\[
-\frac{\beta}{\qxp_1}
-\frac{\eta}{\qxp_1}
+\frac{2\eta}{\qxp}
-\frac{\eta}{2}
\leq0.
\]
Then there exist constants $C_1>0$ and $m_0\in\N$ such that, for every
$m\in\N$ with $m\geq m_0$,
\[
\mathcal{B}_{\Gamma,\qxp,\qxp_1}^m
\leq
C_1
m^{
-\frac{\beta}{\qxp_1}
-\frac{\eta}{\qxp_1}
+\frac{2\eta}{\qxp}
-\frac{\eta}{2}
}.
\]

\end{example}

%%
%%
%%
%%
%%


\begin{example}\label{ex:polynomial-alpha-horn-domain}
Let \(\alpha>0\) and \(m_0\geq2\). Let \(\omb\subset\R^2\) be a smooth domain, and endow \(\mathbb R^2\) with the Euclidean metric \(g=\delta_{\R^2}\).
Assume that
\[
 \overline{\omb}\cap\{(X,Y)\in\R^{2} \mid X\leq m_{0}\}
\]
is compact and that
\[
 \omb\cap\{X>m_{0}\}
 =
 \left\{
 (X,Y)\in\R^{2} \mid
 X>m_{0},\ |Y|<\frac{1}{X^\alpha}
 \right\}.
\]

 
 

Let
\[
 1\leq \qxp,\qxpz,\qxpo<\infty,
 \qquad
 \eta,\beta,\delta,\gamma\in\R.
\]
Consider the weights
\[
 \wghtv_1(p):=(1+|p|)^{1-\beta},
 \qquad
 \wghtv_0(p):=(1+|p|)^{\qxp-\eta},
\]
and
\[
 \wghtv_2(p):=(1+|p|)^\delta,
 \qquad
 \wghtw(p):=(1+|p|)^\gamma.
\]
Assume
\[
 \beta\geq1, \qquad \eta \geq  q,  
 \qquad
 \beta\geq \eta+( \alpha-1)\qxp+1.
\]

Fix
\[
 r\in
 \left(
 \max\left\{\frac12,\frac1{\alpha+1}\right\},
 1
 \right).
\]
Choose \(N\in\N\) such that
\[
N\geq m_0^{\alpha+1}.
\]
{ 
For \(i\in\N\), set
\[
a_i:=(N+i)^{1/(\alpha+1)}.
\]
For this example, take \(D_m=B(0,m)\).
}

Set
\[
r_{0,i}=r_{1,i}:=r,
\qquad
\hat r_{0,i}=\hat r_{1,i}:=1.
\]

For \((s,t)\in Q(0,1)=(-1,1)^2\), define
\[
 X_i(s):=a_i+\frac{s}{a_i^\alpha}.
\]
The interior charts on the tail are
\[
 \varphi_{0,i}(s,t)
 :=
 \left(
 X_i(s),
 \frac{t}{X_i(s)^\alpha}
 \right),
 \qquad i\in\N.
\]
The upper and lower boundary charts on the tail are
\[
 \psi_{1,i}^{\pm}(s,t)
 :=
 \left(
 X_i(s),
 \pm\frac{1-t}{X_i(s)^\alpha}
 \right),
 \qquad i\in\N.
\]

After adjoining finitely many ordinary interior charts and finitely many
ordinary boundary charts on the compact part, and after relabeling the
families
\[
 \varphi_{0,i},
 \qquad
 \psi_{1,i}^{+},
 \qquad
 \psi_{1,i}^{-}
\]
as a family \(\psi_{k,j}\), \((k,j)\in\{0,1\}\times\N\), the resulting
covering satisfies conditions \ref{5}--\ref{3},
with cubes replacing balls as permitted by Remark \ref{109}.

Moreover, there exists a constant \(C>0\), independent of \(i\), such that
\begin{gather*}
 \frac1{Ca_i^\alpha}
 \leq
 \bnorm d\varphi_{0,i}\bnorm, \ \bnorm d\psi_{1,i}^{\pm}\bnorm
 \leq
 \frac{C}{a_i^\alpha}, \qquad
 \frac{a_i^\alpha}{C}
 \leq
 \bnorm d\varphi_{0,i}^{-1}\bnorm, \ \bnorm d(\psi_{1,i}^{\pm})^{-1}\bnorm
 \leq
 Ca_i^\alpha, \\[3pt]
 \frac{1}{Ca_i^{2\alpha}}
 \leq
 \bnorm G_{\varphi_{0,i}}\bnorm , \ \bnorm G_{\psi_{1,i}^{\pm}}\bnorm
 \leq
 \frac{C}{a_i^{2\alpha}},\qquad
 \frac{a_i^{2\alpha}}{C}
 \leq
  \bnorm G_{\varphi_{0,i}}^{-1}\bnorm, \ \bnorm G_{\psi_{1,i}^{\pm}}^{-1}\bnorm
 \leq
 Ca_i^{2\alpha},\\[3pt]
 \frac1{Ca_i^\alpha}
 \leq
 \bnorm G_{\Gamma,\psi_{1,i}^{\pm}}\bnorm
 \leq
 \frac{C}{a_i^\alpha},\\[3pt]
 \kp_{\varphi_{0,i}}
 \leq
 C_r a_i^\alpha,
 \qquad
 \kp_{\psi_{1,i}^{\pm}}
 \leq
 C_r a_i^\alpha.
 \end{gather*}
 
There exist positive continuous functions $b_1$, $b_2$, and $b_3$ such that, on the tail,
\[
b_1(s,t)\approx (1+|(s,t)|)^{1-\beta},
\qquad
b_2(s,t)\approx(1+|(s,t)|)^{\delta},
\qquad
b_3(s,t)\approx(1+|(s,t)|)^{\gamma}.
\]
With this choice, the weights  satisfy conditions \ref{138}--\ref{16}.

If
\[
 \frac{\delta-2\alpha}{\qxpz}
 +
 \frac{\beta+2\alpha-1}{\qxp}
 \leq0,
\]
then there exist constants \(C>0\) and \(m_\ast\geq1\), independent of
\(m\), such that, for every \(m\geq m_\ast\),
\[
 \Bint
 \leq
 C
 m^{
 \frac{\delta-2\alpha}{\qxpz}
 +
 \frac{\beta+2\alpha-1}{\qxp}
 }.
\]
If
\[
 \frac{\gamma-\alpha}{\qxpo}
 +
 \frac{\beta+2\alpha-1}{\qxp}
 \leq0,
\]
then there exist constants \(C>0\) and \(m_\ast\geq1\), independent of
\(m\), such that, for every \(m\geq m_\ast\),
\[
 \Bbd
 \leq
 C
 m^{
 \frac{\gamma-\alpha}{\qxpo}
 +
 \frac{\beta+2\alpha-1}{\qxp}
 }.
\]

\end{example}



\subsubsection{{Analysis of Embeddings}} Our first main results are Propositions \ref{33} and \ref{56}. We write
\[
\Omega_m:=D_m\cap\Omega,
\qquad
\Omega^m:=\Omega\setminus\overline{\Omega_m},
\qquad
\Gamma_m:=D_m\cap\Gamma,
\qquad
\Gamma^m:=\Gamma\setminus\overline{\Gamma_m}.
\]
\begin{proposition} \label{33}
Assume that conditions \ref{5}--\ref{3}, \ref{139},  and \ref{46} hold. Suppose that \(1\leq\qxp<n\) and
$$
\frac{n\qxp }{ n - \qxp } \geq \qxp _0 \geq \qxp.
$$
Then the following statements hold:
\begin{enumerate}[label=(\roman*)]

\item If 
$$
\lim _{m \rightarrow \infty} \mathcal{B} _{ \qxp , \qxp _0 } ^m  <\infty,
$$
then 
$$
W^{1, \qxp  }(\Omega ;  \wghtv _0, \wghtv _1 ) \rightarrow L^{ \qxp _0 }  ( \Omega ; \wghtv _2 )$$ 
is continuous.

\item If  
$$
\qxp _0<\frac{n\qxp}{n-\qxp} \quad \text { and } \quad \lim _{m \rightarrow \infty} \mathcal{B} _{ \qxp , \qxp _0 } ^m =0,
$$ 
then 
$$
W^{1, \qxp  } (\Omega ; \wghtv _0,  \wghtv _1 ) \rightarrow L^{ \qxp _0 }  ( \Omega ; \wghtv _2  )$$
 is compact.
\end{enumerate}
\end{proposition}











\begin{proposition} \label{56}
Assume that conditions \ref{5}--\ref{3}, \ref{139},  and \ref{16} hold. Suppose that \(1\leq\qxp<n\) and
\[
 \frac{(n-1)\qxp }{n - \qxp}  \geq \qxp _1 \geq \qxp.
\]
Then the following statements hold:
 \begin{enumerate}[label=(\roman*)]
\item If  
$$
\lim _{m \rightarrow \infty} \mathcal{B} _{ \Gamma , \qxp , \qxp _1 } ^m <\infty,
$$
 then there exists a continuous trace operator $$
 W^{1, \qxp  } ( \Omega ; \wghtv_0, \wghtv_1 ) \rightarrow L^{\qxp _1}  (\Gamma ; \wghtw).
 $$

\item  If 
$$
\qxp _1<\frac{(n-1)\qxp}{n-\qxp} \quad \text { and } \quad \lim _{m \rightarrow \infty} \mathcal{B} _{ \Gamma , \qxp , \qxp _1 } ^m =0,
$$
 then the trace operator 
 $$
 W^{1, \qxp  } (\Omega ; \wghtv _0, \wghtv _1 ) \rightarrow L^{\qxp _1 }  (\Gamma ; \wghtw)
 $$
  is compact.
\end{enumerate}
\end{proposition}

%In Corollaries \ref{150} and \ref{151}, we additionally obtain results analogous to Propositions \ref{33} and \ref{56} for the case of weighted Sobolev spaces \(W^{1, \qxp}(TM; \wghtv_0, \wghtv_1)\) for the tangent bundle.

\begin{corollary}\label{57}
%Let $\qxp\in[1,n)\to\qxv\in\mathbb R^+ $ and $\qxp\in[1,n)\to\qxw\in\mathbb R^+$ be functions satisfying\[ \qxp<\qxv\leq\frac{n\qxp}{n-\qxp}\qquad\text{and}\qquad\qxp<\qxw\leq\frac{(n-1)\qxp}{n-\qxp}\]for every \(1\leq\qxp<n\).

Suppose that \(n\geq3\) and that conditions \ref{5}--\ref{3} and \ref{139}--\ref{16} hold for \(\qxp=2\). Assume further that
\[
\lim_{m\rightarrow\infty}\mathcal B_{2,\xv}^m<\infty
\quad\text{and}\quad
\lim_{m\rightarrow\infty}\mathcal B_{\Gamma , 2,\xw}^m<\infty.
\]
Then there exist constants \(\mathcal C_{\xv}>0\) and \(\mathcal C_{\xw}>0\) such that
\begin{gather}
\|u\|_{\xv,\Omega,\wghtv_2}
\leq
\mathcal C_{\xv}\|u\|_{1,2,\Omega,\wghtv_0,\wghtv_1},
\label{62}\\[5pt]
\|u\|_{\xw,\Gamma,\wghtw}
\leq
\mathcal C_{\xw}\|u\|_{1,2,\Omega,\wghtv_0,\wghtv_1},
\label{63}
\end{gather}
for every \(u\in W^{1,2}(\Omega;\wghtv_0,\wghtv_1)\).
\end{corollary}

%{\color{red}In the study of problem \eqref{61}, we use \(\xv:=\qxv(2)\) and \(\xw:=\qxw(2)\).}
%Let (\Omega = \mathbb{R}^n \setminus B_r(0)), where (n \geq 3), 









%For our next main result, we look for weak solution $u\in W^{1,2} (\Omega ; \wghtv_0 , \wghtv _1)$ such that the extension of $u$,  see Proposition \ref{22}, belong to $W^{1,2} _0 (M; \wghtv _0 , \wghtv _1 )$, where $ W^{1,2} _0 (M; \wghtv _0 , \wghtv _1 )$ is the completion of $C^{\infty} _0(M)$ with respect to the norm $\| \cdot \| _{ 1 , 2 , M , \wghtv _0 ,  \wghtv _1  } $.


\subsubsection{Study of Eigenvalues} We now assume the following condition.
\begin{enumerate}[label=($\mathtt{W}_{\arabic*}$)]
\setcounter{enumi}{3}
\item \label{47} For every \((k,i)\in\{0,1\}\times\mathbb N\) and almost every \(\pvr\in\hat U_{k,i}\),
\[
{
\kp_{k,i}^{2}(\pvr)\wghtv_1(\pvr)
\leq\ck_2\wghtv_2(\pvr)
}
\quad\text{and}\quad
\int_{\hat U_{k,i}}\wghtv_2\,\dvg\leq\ck_3.
\]
\end{enumerate}


%For weights $V$ and $W$, we adopt the simplified notation

For weights \(\wghtv\) and \(\wghtw\), we use the notation
\[
\dwv=\wghtv\,\dvg,
\qquad
\dww=\wghtw\,\dsg,
\qquad
\wghtv(U)=\int_U\wghtv\,\dvg,
\qquad
\wghtw(U\cap\Gamma)=\int_{U\cap\Gamma}\wghtw\,\dsg,
\]
for every measurable set \(U\subset M\).

%Assume the hypotheses of Corollary 1.4 hold. Suppose $V_1 \in C^1(M)$, $\tau \in L^{2_V / (2_V - 2)}(\Omega; V_2) \cap L^{\infty}(\Omega)$, and the sets $\{\tau > 0\}$ and $\{\tau < 0\}$ have nonempty interiors. Then there exist infinitely many eigenvalues

\begin{theorem}\label{59}
Assume the hypotheses of Corollary \ref{57}. Suppose that \(\wghtv_1\in C^1(M)\),
\[
\iwe\in L^{\xv/(\xv-2)}(\Omega;\wghtv_2)\cap L^\infty(\Omega),
\]
and that the sets
\[
\{\pvr\in\Omega\mid\iwe(\pvr)>0\}
\quad\text{and}\quad
\{\pvr\in\Omega\mid\iwe(\pvr)<0\}
\]
have nonempty interiors. Then problem \eqref{61} has two infinite sequences of eigenvalues
\[
\cdots\leq\lambda_2^-\leq\lambda_1^-<0<\lambda_1^+\leq\lambda_2^+\leq\cdots,
\]
{with \(\lambda_k^-\rightarrow-\infty\) and \(\lambda_k^+\rightarrow+\infty\) as \(k\rightarrow\infty\).}

If \(u\in\womd\) is an eigenfunction of \eqref{61} and condition \ref{47} holds, then \(u\in L^\infty(\Omega)\). Moreover, if \(u\in W_0^{1,2}(\overline\Omega;\wghtv_0,\wghtv_1)\), then \eqref{87} holds.

{Here, \(W_0^{1,2}(\overline\Omega;\wghtv_0,\wghtv_1)\) denotes the closure in \(W^{1,2}(\Omega;\wghtv_0,\wghtv_1)\) of the restrictions to \(\Omega\) of functions in \(C_0^\infty(M)\).}
\end{theorem}

%For some conditions that guarantee \( u \in W^{2,2} (\Omega; \wghtv_0, \wghtv_1, \wghtv_3) \), see Proposition \ref{136} and Remark \ref{156} in Section \ref{143}.


%


%This paper is organized as follows. In Section 2, we collect preliminary definitions and results that are used several times throughout the paper. In Section 3, we study the operators

\medskip

The remainder of the paper is organized as follows. Section \ref{141} presents the preliminary definitions and results. In Section \ref{142}, we establish an extension result and prove continuous and compact weighted Sobolev embedding and trace theorems. Section \ref{143} is devoted to elliptic estimates and the analysis of problem \eqref{61}. In particular, we prove the existence of infinitely many positive and negative eigenvalues and, under additional assumptions, establish the boundedness and decay at infinity of the corresponding eigenfunctions. Finally, the appendix contains the proof of a technical lemma used to derive the elliptic estimates.
  
   
  % In Section \ref{144}, using the spectrum associated with problem \eqref{61}, we establish an existence result for the linear eigenvalue problem \eqref{66} with an indefinite weight and data in \(L^1\). Our approach is non-variational and relies on the notion of solutions obtained by approximation.


  %  In Section \ref{145}, our objective is to compare the first eigenvalue $\lambda _1 (\Omega)$ of \eqref{61}, in the case where $\iwe |_{\Omega}\geq \delta >0$, with the eigenvalue $\lambda _1 ^\pm (M)$ of \eqref{146}. For this purpose, we introduce a parameter motivated by the Dirichlet capacity.  
    
    
    
%%%%%%%%%%%%%%%%%%%%%%%%%%%%%%%%%%%%%%%%%%%%%%%%%%%%%%%%%%%%%%%%%%%%%%%%%%%%%%%%%%%%%%%%%%%%%%%%%%%%%%%%%%%%%%%%%%%%%%%%%%%%%%%%%%%%%%%%%%%%%%%%%%%%%%%%%%%%%%%%%%%%%%%%%%%%%%%%%%%%%%%%%%%%%%%%%%%%%%%%%%%%%%%%%%%%%%%%%%%%%%%%%%%%%%%%%%%%%%%%%%%%%%%%%%%%%%%%%%%%%%%%%%%%%%%%%%%%%%%%%%%%%%%%%%%%%%%%%%%%%%%%%%%%%%%%%%%%%%%%%%%%%%%%%%%%%%%%%%%%%%%%%%%%%%%%%%%%%%%%%%%%%%%%%%%%%%%%%%%%%%%%%%%%%%%%%%%%%%%%%%%%%%%%%%%%%%%%%%%%%%%%%%%%%%%%%%%%%%%%%%%%%%%%%%%%%%%%%%%%%%%%%%%%%%%%%%%%%%%%%%%%%%%%%%%%%%%%%%%%%%%%%%%%%%%%%%%%%%%%%%%%%%%%%%%%%%%%%%%%%%%%%%%%%%%%%%%%%%%%%%%%%%%%%%%%%%%%%%%%%%%%%%%%%%%%%%%%%%%%%%%%%%%%%%%%%%%%%%%%%%%%%%%%%%%%%%%%%%%%%%%%%%%%%%%%%%%%%%%%
    
%%%%%%%%%%%%%%%%%%%%%%%%%%%%%%%%%%%%%%%%%%%%%%%%%%%%%%%%%%%%%%%%%%%%%%%%%%%%%%%%%%%%%%%%%%%%%%%%%%%%%%%%%%%%%%%%%%%%%%%%%%%%%%%%%%%%%%%%%%%%%%%%%%%%%%%%%%%%%%%%%%%%%%%%%%%%%%%%%%%%%%%%%%%%%%%%%%%%%%%%%%%%%%%%%%%%%%%%%%%%%%%%%%%%%%%%%%%%%%%%%%%%%%%%%%%%%%%%%%%%%%%%%%%%%%%%%%%%%%%%%%%%%%%%%%%%%%%%%%%%%%%%%%%%%%%%%%%%%%%%%%%%%%%%%%%%%%%%%%%%%%%%%%%%%%%%%%%%%%%%%%%%%%%%%%%%%%%%%%%%%%%%%%%%%%%%%%%%%%%%%%%%%%%%%%%%%%%%%%%%%%%%%%%%%%%%%%%%%%%%%%%%%%%%%%%%%%%%%%%%%%%%%%%%%%%%%%%%%%%%%%%%%%%%%%%%%%%%%%%%%%%%%%%%%%%%%%%%%%%%%%%%%%%%%%%%%%%%%%%%%%%%%%%%%%%%%%%%%%%%%%%%%%%%%%%%%%%%%%%%%%%%%%%%%%%%%%%%%%%%%%%%%%%%%%%%%%%%%%%%%%%%%%%%%%%%%%%%%%%%%%%%%%%%%%%%%%%%%%%%

%%
%%
%%
%%
%%
%%

\section{Preliminaries} \label{141}

Let $(M,g)$ be a smooth Riemannian manifold. For every { nonnegative} integer $k$ and every smooth function $u:M\rightarrow\mathbb{R}$, we denote by $\nabla^k u$ the $k$th covariant derivative of $u$. {  We set $\nabla^0u=u$.} For $k\geq1$, the pointwise norm of $\nabla^k u$ is defined in local coordinates by
\[
{ 
\left|\nabla^k u\right|^2
=
g^{i_1j_1}\cdots g^{i_kj_k}
(\nabla^k u)_{i_1\ldots i_k}
(\nabla^k u)_{j_1\ldots j_k}.
}
\]

Recall that $(\nabla u)_i=\partial_i u$, while
\begin{equation}\label{137}
{ 
(\nabla^2u)_{ij}=\partial_{ij}u-\Gamma_{ij}^{\ell}\partial_{\ell}u.
}
\end{equation}

We use the following weighted Sobolev spaces; see \cite{kilpel1994weightedcapacity, turesson2000potential, hebey2000nonlinearanaly, kristaly2010variationalprinciples} for related definitions in Euclidean and Riemannian settings. Let $U$ be an open subset of $M$ and let $\qxp\geq1$. For $i=0,1,2$, define
\[
L^{\qxp}(U;\wghtv_i)
:=
\left\{
 u:U\rightarrow\mathbb{R}
 \mid
 u\text{ is measurable and }
 \int_U|u|^{\qxp}\,\dwv_i<\infty
\right\}.
\]

We denote
\begin{gather*}
{ 
\mathcal{C}^{k,\qxp}(\Omega;\wghtv_0,\ldots,\wghtv_k)
:=
\left\{
 u\in C^k(\Omega)
 \mid
 \sum_{i=0}^k\int_\Omega|\nabla^iu|^{\qxp}\,\dwv_i<\infty
\right\},
\qquad k=1,2,
}
\\[5pt]
{ 
C_0^\infty(\Omega)
:=
\left\{
 u\in C^\infty(\Omega)
 \mid
 \spt u\Subset\Omega
\right\},
}
\\[5pt]
C_0^\infty(\overline\Omega)
:=
\left\{
 u\in C^\infty(\Omega)
 \mid
 \text{there exists }\overline u\in C_0^\infty(M)
 \text{ such that }u=\overline u|_\Omega
\right\},
\\[5pt]
{ 
L^{\qxp}_{\loc}(\overline\Omega)
:=
\left\{
 u:\Omega\rightarrow\mathbb{R}
 \mid
 u\in L^{\qxp}(U\cap\Omega)
 \text{ for every open set }U\Subset M
\right\}.
}
\end{gather*}
Recall that $\wghtv_0$ and $\wghtv_1$ are bounded from above on compact subsets of $M$.

\begin{definition}$ $
\begin{enumerate}[label=(\roman*)]
%\item The Sobolev space $W^{k,\qxp}(M;\wghtv_0,\ldots,\wghtv_k)$, $k=1,2$, is the completion of $\mathcal{C}^{k,\qxp}(M;\wghtv_0,\ldots,\wghtv_k)$ with respect to the norm
%\[\|u\|_{k,\qxp,M,\wghtv_0,\ldots,\wghtv_k}:=\left(\sum_{i=0}^k\int_M|\nabla^iu|^{\qxp}\,\dwv_i\right)^{\frac1{\qxp}}.\]

{ 
\item The Sobolev space $W^{k,\qxp}(\Omega;\wghtv_0,\ldots,\wghtv_k)$, $k=1,2$, is the completion of $\mathcal{C}^{k,\qxp}(\Omega;\wghtv_0,\ldots,\wghtv_k)$ with respect to the norm
\[
\|u\|_{k,\qxp,\Omega,\wghtv_0,\ldots,\wghtv_k}
:=
\left(
\sum_{i=0}^k\int_\Omega|\nabla^iu|^{\qxp}\,\dwv_i
\right)^{\frac1{\qxp}}.
\]
}

{ 
\item The Sobolev space $W_0^{k,\qxp}(\overline\Omega;\wghtv_0,\ldots,\wghtv_k)$, $k=1,2$, is the completion of $C_0^\infty(\overline\Omega)$ with respect to the norm $\|\cdot\|_{k,\qxp,\Omega,\wghtv_0,\ldots,\wghtv_k}$.
}

{ 
\item The Sobolev space $W_0^{k,\qxp}(\Omega;\wghtv_0,\ldots,\wghtv_k)$, $k=1,2$, is the completion of $C_0^\infty(\Omega)$ with respect to the norm $\|\cdot\|_{k,\qxp,\Omega,\wghtv_0,\ldots,\wghtv_k}$.
}
\end{enumerate}
\end{definition}

\begin{definition}
We say that $u\in\womd$ is a weak solution of \eqref{61} if
\[
\int_\Omega g(\nabla u,\nabla v)\,\dwv_1
+
\int_\Omega uv\,\dwv_0
=
\lambda\int_\Omega\iwe uv\,\dwv_2
\]
for every $v\in\womd$.
\end{definition}

%%%%
%%%%
%%%%
%%%%
%%%%
%%%%
%%%%



\section{Embeddings of Weighted Sobolev Spaces} \label{142}

%%
%%

In this section, we study continuous and compact embeddings and traces
for the weighted Sobolev space $W^{1,\qxp}(\Omega;\wghtv_0,\wghtv_1)$.

The proofs combine local Euclidean inequalities with the geometry of
the charts and the behavior of the weights. The quantities
\(\mathcal B_{\qxp,\qxp_0}^m\) and
\(\mathcal B_{\Gamma,\qxp,\qxp_1}^m\) control the behavior of functions
in the noncompact part of \(\Omega\). The section begins with an
extension result, followed by the embedding and trace theorems.

 

\begin{lemma} \label{6}
Let $(k,i)\in \{0,1\}\times \mathbb{N}$ and $\delta \geq 2$. There exists a smooth function $\zeta : \hat U_{k,i} \rightarrow [0,1]$ such that
$$
\zeta = 1 \text { in } U_{k,i}, \quad \spt \zeta \subset \psi_{k,i}\left(B\left(0,r_{k,i}+\delta^{-1}(\hat r_{k,i}-r_{k,i})\right)\right),
$$
and
$$
|\nabla \zeta|_g \leq C(n)\bnorm d(\psi_{k,i}^{-1})\bnorm\delta(\hat r_{k,i}-r_{k,i})^{-1} \quad \text { in } \hat U_{k,i}.
$$
\end{lemma}

\begin{proof}
There exists a smooth function $\zeta_0:B(0,\hat r_{k,i})\subset\mathbb R^n\rightarrow[0,1]$ such that
$$
\zeta_0=1 \text { in } B(0,r_{k,i}), \quad \spt\zeta_0\subset B\left(0,r_{k,i}+\delta^{-1}(\hat r_{k,i}-r_{k,i})\right),
$$
and
$$
|\nabla\zeta_0|_{\mec}\leq C_1(n)\delta(\hat r_{k,i}-r_{k,i})^{-1}.
$$
Define $\zeta:=\zeta_0\circ\psi_{k,i}^{-1}$. Then
$$
\zeta=1 \text { in } U_{k,i}, \quad \spt\zeta\subset\psi_{k,i}\left(B\left(0,r_{k,i}+\delta^{-1}(\hat r_{k,i}-r_{k,i})\right)\right),
$$
and
\begin{equation*}
\begin{aligned}
|\nabla\zeta|_g
\leq &\,\displaystyle \bnorm d(\psi_{k,i}^{-1})\bnorm\cdot|\nabla\zeta_0|_{\mec}\
\leq &\,\displaystyle C_1(n)\bnorm d(\psi_{k,i}^{-1})\bnorm\delta(\hat r_{k,i}-r_{k,i})^{-1}.
\end{aligned}
\end{equation*}
This proves the lemma.
\end{proof}

\begin{proposition} \label{22}
Assume that conditions \ref{5}--\ref{3} hold for $\qxp\geq1$, and that the weight functions $\wghtv_0$ and $\wghtv_1$ satisfy conditions \ref{138}--\ref{139}. Let $\delta\geq2$. Then there exists a bounded linear \textnormal{extension operator}
$$
\mathcal E:W^{1,\qxp}(\Omega;\wghtv_0,\wghtv_1)\rightarrow W^{1,\qxp}(M;\wghtv_0,\wghtv_1)
$$
such that
\begin{gather*}
\mathcal Eu=u \text { a.e. in } \Omega,\\[5pt]
\spt\mathcal Eu\subset\bigcup_{(k,i)\in\{0,1\}\times\mathbb N}\psi_{k,i}\left(B\left(0,r_{k,i}+\delta^{-1}(\hat r_{k,i}-r_{k,i})\right)\right),
\end{gather*}
and
\begin{equation}\label{9}
\begin{aligned}
\|\mathcal Eu&\|_{1,\qxp,M,\wghtv_0,\wghtv_1}\\[3pt]
\leq &\,\displaystyle C(n,\qxp)\rd_1^{\frac{q-1}{q}}\left[(1+\ck_0\rd_2)(1+\delta^{\qxp}\ck_{\qxp})\int_\Omega|u|^{\qxp}\,\dwv_0
+\left(1+\ck_0\rd_2^{1+\qxp}\right)\int_\Omega|\nabla u|^{\qxp}\,\dwv_1\right]^{\frac1\qxp}.
\end{aligned}
\end{equation}
\end{proposition}

\begin{proof}

Using Lemma \ref{6}, there exists a partition of unity $\{\zeta _{k,i} : M \rightarrow [0,1]\} _{(k,i) \in \{0,1\} \times \mathbb{N}} \subset C^{\infty}(M)$   subordinate to $\{ \hat{U}_{k,i} \}_{ (k,i) \in \{0,1\} \times \mathbb{N} }$ such that  
$$
\spt \zeta _{k,i} \subset \psi _{k,i} \left( B\left(0 ,  r _{k,i} + \delta ^{-1}(\hat r _{k,i} -r _{k,i} ) \right) \right),    \quad \sum_{k,i} \zeta _{k,i}=1 \text { on }\Omega,
$$ 
and 
$$ 
|\nabla \zeta _{k,i} |_g \leq  C_1(n) \delta \kp _{k , i} \quad \text {  in } \quad \hat U _{k,i}.
$$


%Since the family is locally finite, only finitely many indices satisfy $\psi_{k,i}(0)\in\overline{D_{m_\ast}}$. By the local bounds on $\wghtv_0$ and $\wghtv_1$, we may increase $\ck_{\qxp}$, if necessary, so that the first inequality in \ref{45} also holds on $\spt\zeta_{k,i}$ for these indices.

  
%%%cuentas  {\color{red}For every $(k,i)\in\{0,1\}\times\mathbb N$, let $\eta_{k,i}$ be the function given by Lemma \ref{6}, and set $S:=\sum_{k,i}\eta_{k,i}.$$ Since the family is locally finite, $S$ is smooth. Moreover, $S\geq1$ on $\overline\Omega$. Choose $h\in C^\infty(\mathbb R)$ such that $0\leq h\leq1$, $h=0$ in $(-\infty,1/2]$, and $h=1$ in $[1,\infty)$. Define $$\zeta_{k,i}:=\left\{\begin{aligned}&h(S)\frac{\eta_{k,i}}{S} && \text {if } S>0,\\&0 && \text { if } S=0.  \end{aligned}\right.$$  Then $\{\zeta_{k,i}\}_{(k,i)\in\{0,1\}\times\mathbb N}\subset C^\infty(M)$, $$\spt\zeta_{k,i}\subset\psi_{k,i}\left(B\left(0,r_{k,i}+\delta^{-1}(\hat r_{k,i}-r_{k,i})\right)\right), \quad \sum_{k,i}\zeta_{k,i}=1 \text { on } \Omega, $$ and $$|\nabla\zeta_{k,i}|_g\leq C_1(n)\delta\kp_{k,i}\quad\text { in }\hat U_{k,i}.$$}



Define $f:\mathbb R^n\rightarrow\mathbb R^n$ by $f(x)=(x_1,x_2,\ldots,|x_n|)$. Let $u\in W^{1,\qxp}(\Omega;\wghtv_0,\wghtv_1)$. For $(k,i)\in\{0,1\}\times\mathbb N$, define
\begin{gather*}
\overline u_{0,i}:=u\zeta_{0,i}\quad\text { in }\hat U_{0,i},\\[3pt]
\overline u_{1,j}(\pvr):=\left(u\circ\psi_{1,j}\circ f\circ\psi_{1,j}^{-1}\right)(\pvr)\zeta_{1,j}(\pvr)
\quad\text { if }\pvr\in\hat U_{1,j}.
\end{gather*}

By \ref{138} and \ref{3},
\begin{align}
\int_M|\overline u_{1,j}|^{\qxp}\,\dwv_0
=&\,\displaystyle\int_{\psi_{1,j}(B(0,\hat r_{1,j})\cap\{x_n>0\})}|\overline u_{1,j}|^{\qxp}\,\dwv_0
+\int_{\psi_{1,j}(B(0,\hat r_{1,j})\cap\{x_n<0\})}|\overline u_{1,j}|^{\qxp}\,\dwv_0\nonumber\\[3pt]
\leq&\,\displaystyle(1+\ck_0\rd_2)\int_{\hat U_{1,j}\cap\Omega}|u|^{\qxp}\,\dwv_0.\label{7}
\end{align}

Furthermore, from $|a+b|^{\qxp}\leq2^{\qxp-1}(|a|^{\qxp}+|b|^{\qxp})$, together with  \ref{3} and \ref{138}--\ref{139},
\begin{align}
\int_{M}|\nabla \overline u _{1,j}|^\qxp \, \dwv _1 
= & \displaystyle \, \int _{\psi _{1,j} (B(0,\hat r  _{1,j}) \cap \{x_n >0\})} |\nabla \overline u _{1,j}|^\qxp \, \dwv _1 + \int _{\psi _{1,j} (B(0,\hat r  _{1,j}) \cap \{x_n < 0\})} |\nabla \overline u _{1,j}|^\qxp \, \dwv _1 \nonumber\\[3pt]
 \leq & \displaystyle \,  \int _{\psi _{1,j} (B(0,\hat r  _{1,j}) \cap \{x_n >0\})} 2^{\qxp-1} \left( | \zeta _{1,j} \nabla u|^\qxp  + |u \nabla \zeta _{1,j} |^\qxp \right) \, \dwv _1  \nonumber \\[3pt]
& \displaystyle \, + \int _{\psi _{1,j} (B(0,\hat r  _{1,j}) \cap \{x_n > 0\})} 2^{\qxp-1} \ck _0 \rd _2 \left( |\rd _2 \zeta _{1,j} \nabla u |^\qxp  + |u \nabla \zeta _{1,j} |^\qxp \right) \, \dwv _1 \nonumber\\[5pt]
 \leq & \displaystyle \,  C_2 (n , \qxp ) \left[\delta ^{\qxp} \ck _{\qxp} (1+ \ck _0 \rd _{2})\int_{\hat U _{1,j} \cap \Omega }|u|^\qxp \, \dwv _0 +  \left(1+ \ck _0 \rd _2 ^{1+\qxp} \right)\int_{\hat U _{1,j} \cap \Omega }|\nabla u|^\qxp \, \dwv _1 \right]. \label{8}
\end{align}

{
For the interior charts, the same computation without reflection gives
\begin{align*}
\int_M|\overline u_{0,i}|^{\qxp}\,\dwv_0
&\leq\int_{\hat U_{0,i}}|u|^{\qxp}\,\dwv_0,\\[3pt]
\int_M|\nabla\overline u_{0,i}|^{\qxp}\,\dwv_1
&\leq C_3(n,\qxp)\left[\delta^{\qxp}\ck_{\qxp}\int_{\hat U_{0,i}}|u|^{\qxp}\,\dwv_0
+\int_{\hat U_{0,i}}|\nabla u|^{\qxp}\,\dwv_1\right].
\end{align*}
}

Define
$$
\mathcal Eu:=\sum_{(k,i)\in\{0,1\}\times\mathbb N}\overline u_{k,i}.
$$
Then $\mathcal E$ is linear and $\mathcal Eu=u$ a.e. in $\Omega$. { Since at most $\rd_1$ terms are nonzero at each point,}
$$
{ \left|\sum_{k,i}\overline u_{k,i}\right|^{\qxp}\leq\rd_1^{\qxp-1}\sum_{k,i}|\overline u_{k,i}|^{\qxp}.}
$$
Hence, by \ref{1}, \eqref{7}, and \eqref{8}, we obtain \eqref{9}.
\end{proof}

\subsection{Compact Embeddings}

The proof of our embedding theorems is based on the following lemma.
\begin{lemma} \label{21}
$ $
\begin{enumerate}[label=(\roman*)]
\item \label{18} Assume that
\begin{gather}
W^{1,\qxp}(\Omega_m;\wghtv_0,\wghtv_1)\rightarrow L^{\qxp_0}(\Omega_m;\wghtv_2) \text { is compact for every }m,\label{10}\\[3pt]
\lim_{m\rightarrow\infty}\sup_{\|u\|_{1,\qxp,\Omega,\wghtv_0,\wghtv_1}\leq1}\|u\|_{\qxp_0,\Omega^m,\wghtv_2}=0.\label{11}
\end{gather}
Then
$$
W^{1,\qxp}(\Omega;\wghtv_0,\wghtv_1)\rightarrow L^{\qxp_0}(\Omega;\wghtv_2)
$$
is compact.

On the other hand, if
$$
W^{1,\qxp}(\Omega;\wghtv_0,\wghtv_1)\rightarrow L^{\qxp_0}(\Omega;\wghtv_2)
$$
is compact, then \eqref{11} holds.

\item \label{19} If
\begin{gather}
W^{1,\qxp}(\Omega_m;\wghtv_0,\wghtv_1)\rightarrow L^{\qxp_0}(\Omega_m;\wghtv_2) \text { is continuous for every }m,\label{12}\\[5pt]
\lim_{m\rightarrow\infty}\sup_{\|u\|_{1,\qxp,\Omega,\wghtv_0,\wghtv_1}\leq1}\|u\|_{\qxp_0,\Omega^m,\wghtv_2}<\infty.\label{13}
\end{gather}
Then
$$
W^{1,\qxp}(\Omega;\wghtv_0,\wghtv_1)\rightarrow L^{\qxp_0}(\Omega;\wghtv_2)
$$
is continuous.

If
$$
W^{1,\qxp}(\Omega;\wghtv_0,\wghtv_1)\rightarrow L^{\qxp_0}(\Omega;\wghtv_2)
$$
is continuous, then \eqref{13} holds.
\end{enumerate}
\end{lemma}

\begin{proof}
$\ref{18}$ Notice that \eqref{11} is equivalent to the statement that for every $\epsilon>0$ there exists $m_0$ such that
\begin{equation}\label{17}
\|u\|_{\qxp_0,\Omega,\wghtv_2}^{\qxp_0}
\leq\epsilon\|u\|_{1,\qxp,\Omega,\wghtv_0,\wghtv_1}^{\qxp_0}
+\|u\|_{\qxp_0,\Omega_{m_0},\wghtv_2}^{\qxp_0}
\quad\forall u\in W^{1,\qxp}(\Omega;\wghtv_0,\wghtv_1).
\end{equation}

Let $(u_j)$ be a bounded sequence in $W^{1,\qxp}(\Omega;\wghtv_0,\wghtv_1)$, with
$$
\|u_j\|_{1,\qxp,\Omega,\wghtv_0,\wghtv_1}\leq C_1.
$$
{By \eqref{10} and a diagonal argument, there exists a subsequence, still denoted by $(u_j)$, which is a Cauchy sequence in $L^{\qxp_0}(\Omega_m;\wghtv_2)$ for every $m\in\mathbb N$.}
Let $\epsilon>0$. From \eqref{17}, for sufficiently large $i,j$,
\begin{align*}
\|u_j-u_i\|_{\qxp_0,\Omega,\wghtv_2}^{\qxp_0}
\leq&\,\displaystyle\epsilon\|u_j-u_i\|_{1,\qxp,\Omega,\wghtv_0,\wghtv_1}^{\qxp_0}
+\|u_j-u_i\|_{\qxp_0,\Omega_{m_0},\wghtv_2}^{\qxp_0}\\[3pt]
\leq&\,\displaystyle\left(2^{\qxp_0}C_1^{\qxp_0}+1\right)\epsilon.
\end{align*}
Consequently, $(u_j)$ is a Cauchy sequence in $L^{\qxp_0}(\Omega;\wghtv_2)$. Hence the embedding is compact.

Now, let the embedding be compact and assume that \eqref{17} does not hold. Then there exist $\epsilon>0$ and a sequence $(u_j)$ in $W^{1,\qxp}(\Omega;\wghtv_0,\wghtv_1)$ such that
$$
\|u_j\|_{\qxp_0,\Omega,\wghtv_2}^{\qxp_0}
>\epsilon\|u_j\|_{1,\qxp,\Omega,\wghtv_0,\wghtv_1}^{\qxp_0}
+\|u_j\|_{\qxp_0,\Omega_j,\wghtv_2}^{\qxp_0}.
$$
Write
$$
\tilde u_j:=\frac{u_j}{\|u_j\|_{1,\qxp,\Omega,\wghtv_0,\wghtv_1}}.
$$
Then
\begin{equation}\label{20}
\|\tilde u_j\|_{\qxp_0,\Omega,\wghtv_2}^{\qxp_0}
>\epsilon+\|\tilde u_j\|_{\qxp_0,\Omega_j,\wghtv_2}^{\qxp_0}.
\end{equation}
Since $(\tilde u_j)$ is bounded in $W^{1,\qxp}(\Omega;\wghtv_0,\wghtv_1)$, there is a subsequence converging to $\tilde u$ in $L^{\qxp_0}(\Omega;\wghtv_2)$. Moreover,
$$
\left|\|\tilde u_j\|_{\qxp_0,\Omega_j,\wghtv_2}-\|\tilde u\|_{\qxp_0,\Omega,\wghtv_2}\right|
\leq\|\tilde u_j-\tilde u\|_{\qxp_0,\Omega,\wghtv_2}
+\left|\|\tilde u\|_{\qxp_0,\Omega_j,\wghtv_2}-\|\tilde u\|_{\qxp_0,\Omega,\wghtv_2}\right|.
$$
Taking the limit in \eqref{20}, we obtain
$$
\|\tilde u\|_{\qxp_0,\Omega,\wghtv_2}^{\qxp_0}
\geq\epsilon+\|\tilde u\|_{\qxp_0,\Omega,\wghtv_2}^{\qxp_0},
$$
which is a contradiction.

\medskip

$\ref{19}$ The second part of the lemma can be proved similarly to $\ref{18}$.
\end{proof}

\subsubsection{Proof of Proposition \ref{33}}

We have
$$
\|v\|_{\qxp_0,B(0,1)}\leq C_1\|v\|_{1,\qxp,B(0,1)\cap\{x_n>0\}},
\quad\forall v\in W^{1,\qxp}(B(0,1)\cap\{x_n>0\}),
$$
where $B(0,1)\subset \mathbb{R}^n$ and $C_1:=C_1(\qxp , \qxp _0 , n , B(0,1)) >0 $.

Let $u\in W^{1,\qxp}(\Omega;\wghtv_0,\wghtv_1)$. Then,
\begin{equation*}
\begin{aligned}
\left(\int_{\hat U_{k,i}\cap\Omega}|u|^{\qxp_0}\,\dvg\right)^{\frac1{\qxp_0}}
= &\,\displaystyle \left(\int_{ B(0 , \hat r_{k,i} )\cap \{x_n >0\} } | u \circ \psi _{k,i} |^{ \qxp _0 }  \sqrt{\det [ g_{\gij} ]} \, \dx \right)^{\frac{1}{ \qxp _0 }}\\[3pt]
  \leq & \displaystyle \,  C_1  \bnorm G_{k,i}\bnorm ^{\frac{1}{ \qxp _0 }}  \hat{r}_{k,i}^{\frac{n}{ \qxp _0  }} \left(\int_{ B(0 , 1 ) \cap \{x_n >0\} } | u \circ \psi _{k,i}(\hat r _{k,i} y) |^{ \qxp _0 } \,  \dy \right)^{\frac{1}{   \qxp _0}  }\\[5pt]
\leq&\,\displaystyle C_2  \bnorm G_{k,i}\bnorm^{\frac1{\qxp_0}}
\bnorm G_{k,i}^{-1}\bnorm^{\frac1\qxp}\hat r_{k,i}^{\frac n{\qxp_0}}\\[3pt]
&\,\displaystyle\cdot\left(\hat r_{k,i}^{-n}\int_{\hat U_{k,i}\cap\Omega}|u|^{\qxp}\,\dvg
+\hat r_{k,i}^{-n+\qxp}\bnorm d\psi_{k,i}\bnorm^{\qxp}
\int_{\hat U_{k,i}\cap\Omega}|\nabla u|^{\qxp}\,\dvg\right)^{\frac1\qxp},
\end{aligned}
\end{equation*}
where $C_2:= C_1 n^{1/2}$.

From  \ref{139} and \ref{46},
\begin{equation}\label{155}
\begin{aligned}
\int_{\hat U_{k,i}\cap\Omega}|u|^{\qxp_0}\,\dwv_2 \leq & \, \displaystyle  \left[b_2  ^{\frac{1}{ \qxp _0 }} (\psi _{k,i} (0)) C_2 \bnorm G_{k,i}\bnorm ^{\frac{1}{  \qxp _0 }}  \bnorm G_{k,i} ^{-1}\bnorm ^{\frac{1}{ \qxp}}   \hat r_{k,i}^{\frac{n}{ \qxp _0   }-\frac{n}{\qxp }+1} \right] ^{ \qxp _0 }   \\[3pt]
& \, \displaystyle   \cdot \left(\hat r_{k,i}^{-\qxp  }\int_{\hat U _{k,i}}| u |^{\qxp} \, \dvg + \bnorm d \psi _{k,i}\bnorm ^{\qxp}  \int_{\hat U _{k,i}}|\nabla  u  |^{\qxp}  \, \dvg \right)^{\frac{ \qxp _0 }{ \qxp}  } \\[3pt]
\leq&\,\displaystyle C_2^{\qxp_0}
\left[\frac{b_2^{\frac1{\qxp_0}}(\psi_{k,i}(0))}{b_1^{\frac1\qxp}(\psi_{k,i}(0))}
\bnorm G_{k,i}\bnorm^{\frac1{\qxp_0}}
\bnorm G_{k,i}^{-1}\bnorm^{\frac1\qxp}
\hat r_{k,i}^{\frac n{\qxp_0}-\frac n\qxp+1}\right]^{\qxp_0}\\[3pt]
&\,\displaystyle\cdot\left(\hat r_{k,i}^{-\qxp}\bnorm d\psi_{k,i}\bnorm^{-\qxp}
\int_{\hat U_{k,i}\cap\Omega}|u|^{\qxp}\,\dwv_1
+\int_{\hat U_{k,i}\cap\Omega}|\nabla u|^{\qxp}\,\dwv_1\right)^{\frac{\qxp_0}\qxp}\\[3pt]
\leq&\,\displaystyle C_2^{\qxp_0}(\mathcal B_{\qxp,\qxp_0}^m)^{\qxp_0}
\left(\ck_{\qxp}\int_{\hat U_{k,i}\cap\Omega}|u|^{\qxp}\,\dwv_0
+\int_{\hat U_{k,i}\cap\Omega}|\nabla u|^{\qxp}\,\dwv_1\right)^{\frac{\qxp_0}\qxp},
\end{aligned}
\end{equation}
where condition \ref{139} was used in the last inequality.

Since
$$
\chi_{\Omega}
\leq\sum_{(k,i)\in\{0,1\}\times\mathbb N} \chi_{\hat U_{k,i}\cap \Omega} \leq \rd _1 \chi _{\Omega},
$$
from \ref{1}, \eqref{155}, and
$$
\sum_i|a_i|^\kappa\leq\left(\sum_i|a_i|\right)^\kappa,
\quad \kappa\geq1,
$$
we obtain
\begin{align*}
\|u\|_{\qxp_0,\Omega^{m},\wghtv_2}^{\qxp_0}
\leq&\,\displaystyle C_2 ^{\qxp_0}(\mathcal B_{\qxp,\qxp_0}^m)^{\qxp_0}
\left[\sum_{k,i}\left(\ck_{\qxp}\int_{\hat U_{k,i}\cap\Omega}|u|^{\qxp}\,\dwv_0
+\int_{\hat U_{k,i}\cap\Omega}|\nabla u|^{\qxp}\,\dwv_1\right)\right]^{\frac{\qxp_0}\qxp}\\[3pt]
\leq&\,\displaystyle C_2 ^{\qxp_0}(\mathcal B_{\qxp,\qxp_0}^m)^{\qxp_0}
[\rd_1(\ck_{\qxp}+1)]^{\frac{\qxp_0}\qxp}
\|u\|_{1,\qxp,\Omega,\wghtv_0,\wghtv_1}^{\qxp_0}.
\end{align*}
Thus,
\begin{equation}\label{24}
\|u\|_{\qxp_0,\Omega^{m},\wghtv_2}
\leq C_2 \mathcal B_{\qxp,\qxp_0}^m[\rd_1(\ck_{\qxp}+1)]^{\frac1\qxp}
\|u\|_{1,\qxp,\Omega,\wghtv_0,\wghtv_1}.
\end{equation}

Recall that, on every compact subset of $M$, the weights are bounded above and below by positive constants. Therefore, the local embedding in \eqref{12} is continuous when $\qxp_0\leq n\qxp/(n-\qxp)$, whereas the local embedding in \eqref{10} is compact when $\qxp_0<n\qxp/(n-\qxp)$. The proposition now follows from Lemma \ref{21} and \eqref{24}.


\begin{flushright}
$\square$
\end{flushright}

\subsection{Compact Trace Results}
Proceeding similarly to Lemma \ref{21}, we have the following lemma.
\begin{lemma} \label{25}
$ $
\begin{enumerate}[label=(\roman*)]
\item \label{26} Assume that
\begin{gather}
W^{1,\qxp}(\Omega_m;\wghtv_0,\wghtv_1)\rightarrow L^{\qxp_1}(\Gamma_m;\wghtw) \text { is compact for every }m,\label{27}\\[5pt]
\lim_{m\rightarrow\infty}\sup_{\|u\|_{1,\qxp,\Omega,\wghtv_0,\wghtv_1}\leq1}\|u\|_{\qxp_1,\Gamma^m,\wghtw}=0.\label{28}
\end{gather}
Then
$$
W^{1,\qxp}(\Omega;\wghtv_0,\wghtv_1)\rightarrow L^{\qxp_1}(\Gamma;\wghtw)
$$
is compact.

On the other hand, if the trace operator is compact, then \eqref{28} holds.

\item \label{29} If
\begin{gather}
W^{1,\qxp}(\Omega_m;\wghtv_0,\wghtv_1)\rightarrow L^{\qxp_1}(\Gamma_m;\wghtw) \text { is continuous for every }m,\label{30}\\[5pt]
\lim_{m\rightarrow\infty}\sup_{\|u\|_{1,\qxp,\Omega,\wghtv_0,\wghtv_1}\leq1}\|u\|_{\qxp_1,\Gamma^m,\wghtw}<\infty.\label{31}
\end{gather}
Then
$$
W^{1,\qxp}(\Omega;\wghtv_0,\wghtv_1)\rightarrow L^{\qxp_1}(\Gamma;\wghtw)
$$
is continuous.

If the trace operator is continuous, then \eqref{31} holds.
\end{enumerate}
\end{lemma}


%%
%%
%%
%%
%%
%%
%%%cuentas $$\begin{aligned} \left(\int_{\hat U_{1,j}\cap\Gamma}|Tu|^{\qxp_1}\,\dsg\right)^{\frac1{\qxp_1}}=\left(\int_{B(0,\hat r_{1,j})\cap\{x_n=0\}}\left|\operatorname{Tr}_{1,j}(u\circ\psi_{1,j})\right|^{\qxp_1}\sqrt{\det[g_{\Gamma,\gij}]}\,\ds\right)^{\frac1{\qxp_1}}.\end{aligned}$$





%%
%%

\subsubsection{Proof of Proposition \ref{56}}
 
 
\noindent {\bf Step 1.} We begin by defining the trace operator. For every
\(j\in\mathbb N\), set
\[
A_j:=\hat U_{1,j}\cap\Gamma.
\]
Let
\[
\operatorname{Tr}_{1,j}:
W^{1,\qxp}
\left(
B(0,\hat r_{1,j})\cap\{x_n>0\}
\right)
\longrightarrow
L^{\qxp_1}
\left(
B(0,\hat r_{1,j})\cap\{x_n=0\}
\right)
\]
be the Euclidean trace operator. For
\(u\in W^{1,\qxp}(\Omega;\wghtv_0,\wghtv_1)\), define the local trace
\(T_j u\) on \(A_j\) by
\[
(T_j u)\circ\psi_{1,j}
:=
\operatorname{Tr}_{1,j}(u\circ\psi_{1,j}).
\]

This definition is well posed because the weights are bounded above
and below by positive constants on compact subsets of \(M\). Hence
\(u\circ\psi_{1,j}\) belongs to the corresponding Euclidean Sobolev
space.

If \(A_j\cap A_l\neq\varnothing\), the transition map
\[
\psi_{1,l}^{-1}\circ\psi_{1,j}
\]
is a smooth change of boundary coordinates. Since the Euclidean trace
commutes with smooth changes of coordinates preserving the flat
boundary, we have
\[
T_j u=T_l u
\quad\text{a.e. on }A_j\cap A_l.
\]

To define the global trace, set
\[
E_1:=A_1,
\qquad
E_j:=A_j\setminus\bigcup_{l=1}^{j-1}A_l,
\quad j\geq2.
\]
The sets \(E_j\) are measurable, pairwise disjoint, and
\[
\Gamma=\bigcup_{j\in\mathbb N}E_j.
\]

We define
\[
Tu:=\sum_{j=1}^{\infty}\chi_{E_j}T_j u.
\]


Thus, \(Tu\) is measurable. Moreover, the compatibility of the local
traces gives
\[
Tu=T_j u
\quad\text{a.e. on }A_j
\]
for every \(j\in\mathbb N\). Since the family
\(\{A_j\}_{j\in\mathbb N}\) is countable and covers \(\Gamma\), this
property also shows that \(Tu\) is unique up to equality almost
everywhere on \(\Gamma\).

In addition, if
\[
u\in C(\overline\Omega)\cap
W^{1,\qxp}(\Omega;\wghtv_0,\wghtv_1),
\]
then
\[
Tu=u|_\Gamma
\quad\text{a.e. on }\Gamma.
\]


%This construction of the trace by boundary coordinates is standard;see \cite{zbMATH06237307}.



%%
%%
%%
%%

\medskip

 
\noindent {\bf Step 2.} We now prove the continuity and compactness assertions for the trace operator.

For simplicity, we identify \(Tu\) with \(u\) on \(\Gamma\); that is, we write
\[
u:=Tu
\qquad\text{on }\Gamma.
\]


We have
$$
\|v\|_{\qxp_1,B(0,1)\cap\{x_n=0\}}
\leq c_1\|v\|_{1,\qxp,B(0,1)\cap\{x_n>0\}},
\quad\forall v\in W^{1,\qxp}(B(0,1)\cap\{x_n>0\}),
$$
where $B(0,1)\subset \mathbb{R}^n$ and $c_1 := c_1(\qxp , \qxp _1 , n , B(0,1))>0$.



For $u\in W^{1,\qxp}(\Omega;\wghtv_0,\wghtv_1)$, we obtain
\begin{align*}
\left(\int_{\hat U_{1,j}\cap\Gamma}|u|^{\qxp_1}\,\dsg\right)^{\frac1{\qxp_1}} = & \,  \displaystyle   \left(\int_{ B(0 , \hat r_{1,j} )\cap \{x_n =0\}  } | u \circ \psi _{1,j} |^{\qxp _1}  \sqrt{\det [ g_{\Gamma, \gij}]} \, \ds \right)^{\frac{1}{\qxp _1}}\\[3pt]
  \leq & \, \displaystyle   c _1 \bnorm  G_{\Gamma , 1,j} \bnorm ^{\frac{1}{\qxp _1}}  \hat{r}_{1,j}^{\frac{n-1}{\qxp _1 }} \left(\int_{ B(0 , 1 )\cap \{x_n >0\}} | u \circ \psi _{1,j}(\hat r _{1,j} y) |^{\qxp }   \dy \right)^{\frac{1}{\qxp}} \\[5pt]
\leq&\,\displaystyle c_2 \bnorm G_{\Gamma,1,j}\bnorm^{\frac1{\qxp_1}}
\bnorm G_{1,j}^{-1}\bnorm^{\frac1\qxp}\hat r_{1,j}^{\frac{n-1}{\qxp_1}}\\[3pt]
&\,\displaystyle\cdot\left(\hat r_{1,j}^{-n}\int_{\hat U_{1,j}\cap\Omega}|u|^{\qxp}\,\dvg
+\hat r_{1,j}^{-n+\qxp}\bnorm d\psi_{1,j}\bnorm^{\qxp}
\int_{\hat U_{1,j}\cap\Omega}|\nabla u|^{\qxp}\,\dvg\right)^{\frac1\qxp},
\end{align*}
where $c_2 := c_1 n^{1/2}$.

Proceeding as in the proof of Proposition \ref{33}, it follows from \ref{16} and \ref{139} that
\begin{align*}
\int_{\hat U_{1,j}\cap\Gamma}|u|^{\qxp_1}\,\dww
\leq&\,\displaystyle c_2 ^{\qxp_1}(\mathcal B_{\Gamma,\qxp,\qxp_1}^m)^{\qxp_1}
\left(\ck_{\qxp}\int_{\hat U_{1,j}\cap\Omega}|u|^{\qxp}\,\dwv_0
+\int_{\hat U_{1,j}\cap\Omega}|\nabla u|^{\qxp}\,\dwv_1\right)^{\frac{\qxp_1}\qxp}.
\end{align*}

From \ref{1},
\begin{align*}
\|u\|_{\qxp_1,\Gamma^{m},\wghtw}^{\qxp_1}
\leq&\,\displaystyle c_2 ^{\qxp_1}(\mathcal B_{\Gamma,\qxp,\qxp_1}^m)^{\qxp_1}
\left[\sum_j\left(\ck_{\qxp}\int_{\hat U_{1,j}\cap\Omega}|u|^{\qxp}\,\dwv_0
+\int_{\hat U_{1,j}\cap\Omega}|\nabla u|^{\qxp}\,\dwv_1\right)\right]^{\frac{\qxp_1}\qxp}\\[3pt]
\leq&\,\displaystyle c_2 ^{\qxp_1}(\mathcal B_{\Gamma,\qxp,\qxp_1}^m)^{\qxp_1}
[\rd_1(\ck_{\qxp}+1)]^{\frac{\qxp_1}\qxp}
\|u\|_{1,\qxp,\Omega,\wghtv_0,\wghtv_1}^{\qxp_1}.
\end{align*}
Thus,
\begin{equation}\label{35}
\|u\|_{\qxp_1,\Gamma^{m},\wghtw}
\leq c_2 \mathcal B_{\Gamma,\qxp,\qxp_1}^m[\rd_1(\ck_{\qxp}+1)]^{\frac1\qxp}
\|u\|_{1,\qxp,\Omega,\wghtv_0,\wghtv_1}.
\end{equation}

Recall that the weights are bounded above and below by positive constants on every compact subset of \(M\). Therefore, the local trace operator in \eqref{30} is continuous when $\qxp_1\leq (n-1)\qxp/(n-\qxp)$,
  whereas the local trace operator in \eqref{27} is compact when $\qxp_1<(n-1)\qxp/(n-\qxp)$. The proposition then follows from Lemma \ref{25} and \eqref{35}.


\begin{flushright}
$\square$
\end{flushright}


\section{Principal Eigenvalues with Indefinite Weights}\label{143}


%\subsection{Elliptic estimates} 

Before proving Theorem \ref{59}, we prove the preliminary Propositions \ref{52} and \ref{53}.

\subsection{Elliptic Estimates} 

Similarly to Lemma \ref{6}, we have the following result.
\begin{lemma} \label{153}
Suppose $0<r<R\leq \hat r _{1,j}$ and fix $j\in \mathbb{N}$. There exists a smooth function $\zeta  : \hat U  _{1,j} \rightarrow [0,1]$ such that  
$$
\zeta  = 1 \text { in } \psi _{1,j} (B(0,r)), \quad   \spt \zeta  \subset \psi _{1,j} (B(0,R)),
$$ 
and 
$$
|\nabla \zeta  | _g \leq \cen _1(n)\bnorm d (\psi _{1,j}^{-1}) \bnorm  (R-r)^{-1} \quad \text { in } \quad \hat  U _{1,j}.
$$ 
\end{lemma}



For $h \in \mathbb{R}$ and $t\in (0,\hat r_{1,j})$, we define
\begin{gather*}
\ops [u, v] :=\int  _{\Omega}    g(  \nabla u , \nabla v ) \, \dwv _{1}  + \int  _{\Omega}  u v  \,   \dwv _{0} ,\\[5pt]
\Psi (h ,t) := \sqrt{\int _{ \mathscr{U}(h ,t) } u_{ h }  ^2  \, \dwv _2 + \int _{ \mathscr{U}_{\Gamma }  (h ,t)  } u_{ h } ^2   \, \dww },
\end{gather*}
where $u,v\in \womd  $,
\begin{gather*}
u_{h}  := (u-h )^+,\\[5pt]
\mathscr{U}(h ,t) := \left\{\pvr \in \psi _{1,j} (  B(0,t)\cap \{x_n >0\} ) \mid u(\pvr) > h  \right\} , \\[5pt]
 \mathscr{U}_{ \Gamma } (h ,t) := \left\{ \pvr \in \psi _{1,j} (  B(0,t)\cap \{x_n =0\} ) \mid u(\pvr) > h  \right\}.  
\end{gather*}



The proof of the next lemma will be given in the appendix.
\begin{lemma} \label{alcontb0}
Suppose $\kp _{1 , j} ^{2 } (\pvr) \wghtv _1 (\pvr )  \leq \ck _{2} \wghtv _2 (\pvr )$ for   a.e. $\pvr  \in \hat U _{1,j}$. Assume that   $\fc _2, f \in L^{q_2} (\hat U _{1,j} \cap \Omega  ; \wghtv _2)$,  and  $\fc _3 , f_1 \in L^{q_3} (\hat U _{1,j} \cap \Gamma ; \wghtw)$  for some  $q _2, q_3 > \max \{ \xv / (\xv -2) , \xw / (\xw -2) \}$. Suppose  
\begin{equation}\label{a168}
\begin{aligned}
 \ops [u,v_{h} ] + \int  _{\hat U _{1,j} \cap \Omega } \fc _2 u v_{h}   \, \dwv _{2}   + \int  _{ \hat U _{1,j} \cap \Gamma} \fc _3 u v_{ h } \dww  
 \leq \int  _{\hat U _{1,j} \cap \Omega }    f v_{ h } \, \dwv _{2}  + \int_{\hat U _{1,j} \cap \Gamma }  f_1 v _{ h } \, \dww ,  
\end{aligned}
\end{equation}
for all $h\geq 0 $, where $v _{ h }  = u_{ h }  \zeta ^2 $ and $\zeta$ is given in Lemma \ref{153}.    


There exist  $\epsilon  = \epsilon ( q_2 , q_3 , \xv , \xw ) >0$ and  $\cen _2 = \cen _2 (  n , \mathcal C_{\xv} , \mathcal C_{\xw} , \xv , \xw , \qxp _2 , \qxp _3) >0$  such that if  
\begin{align*}
h_{2} > & \displaystyle \, h_{1}  \\[3pt]
  \geq & \displaystyle \,  \cen _2 \max \left\{ \widehat{\fc}_2 ^{\frac{\qxp _2 \xv }{\qxp _2 (\xv -2) - \xv}}   , \widehat{\fc}_3 ^{\frac{\qxp _3 \xw }{\qxp _3 (\xw -2) - \xw}} , 1\right\}  \max \left\{ \Vert u^+ \Vert _{2 , \hat U _{1,j} \cap \Omega  , \wghtv _2}  ,   \Vert  u^+ \Vert _{2 , \hat U _{1,j} \cap \Gamma , \wghtw }  \right\}, 
\end{align*}
then 
\begin{equation} \label{48}
\begin{aligned}
&\Psi  (h_{2},r) \\[3pt]
&\quad \ \leq \cen _3 \left[   \frac{  \hat r_{1,j}  - r_{1,j} }{(R-r)(h_{2}-h_{1} )^{\epsilon }}  + \frac{ \Vert f \Vert _{ q_2 ,  \hat U _{1,j} \cap \Omega  , \wghtv _2}  + \Vert f_1 \Vert _{q_3 , \hat U _{1,j} \cap \Gamma , \wghtw} + h_{2} (\widehat{\fc}_2 + \widehat{\fc}_3)^{1/2}}{(h_{2}-h_{1} )^{1+\epsilon }}\right] \Psi ^{1 + \epsilon  } (h_{1} ,R),
\end{aligned}
\end{equation}
where  $\cen _3= \cen _3  (  n , \mathcal C_{\xv} , \mathcal C_{\xw} ,  \xv , \xw , \qxp _2 , \qxp _3 , \ck _2) >0$, and
$$
\widehat{\fc}_2 :=\Vert \fc _2 \Vert _{q_2 , \hat U _{1,j} \cap \Omega  , \wghtv _2 } ,\qquad \widehat{\fc}_3 :=\Vert \fc _3 \Vert _{q_3 , \hat U _{1,j} \cap \Gamma  , \wghtw } .
$$

\end{lemma}
%%%%%%%%%%%%%%%%%%%%%%%%%%%%%%%%%%%%%%%%%%%%%%%%%%%%%%%%%%%%%%%%%%%%%%%%%%%%%%%%%%%%%%%%%%%%%%%%%%%%%%%%%%%%%%%%%%%%%%%%%%%%%%%%%%%%%%%%%%%%%%%%%%%%%%%%%%%%%%%%%%%%%%%%%%%%%%

 
%%%%%%%%%%%%%%%%%%%%%%%%%%%%%%%%%%%%%%%%%%%%%%%%%%%%%%%%%%%%%%%%%%%%%%%%%%%%%%%%%%%%%%%%%%%%%%%%%%%%%%%%%%%%%%%%%%%%%%%%%%%%%%%%%%%%%%%%%%%%%%%%%%%%%%%%%%%


\begin{proposition} \label{52}
Suppose $\kp _{1 , j} ^{2 } (\pvr) \wghtv _1 (\pvr )  \leq \ck _{2} \wghtv _2 (\pvr )$ for   a.e. $\pvr  \in \hat U _{1,j}$.  Assume that   $\fc _2, f \in L^{q_2} (\hat U _{1,j} \cap \Omega  ; \wghtv _2)$,  and  $\fc _3 , f_1 \in L^{q_3} (\hat U _{1,j} \cap \Gamma ; \wghtw)$  for some  $q _2, q_3 > \max \{ \xv / (\xv -2) , \xw / (\xw -2) \}$. Suppose  
\begin{equation*}
\begin{aligned}
 \ops [u,v_{h} ] + \int  _{\Omega } \fc _2 u v_{h}   \,  \dwv _{2}   + \int  _{ \Gamma }\fc _3 u v_{ h } \, \dww  
 \leq \int  _{\Omega }    f v_{ h } \,  \dwv _{2}  + \int_{\Gamma } f_1 v _{ h } \, \dww ,  
\end{aligned}
\end{equation*}
for all $h\geq 0 $,   $v _{ h }  = u_{ h }  \zeta ^2 $ and $\zeta \in C^{\infty} _0 ( \hat U_{1,j} )$. Then
\begin{align*}
 \stackbin[ U _{1,j} \cap \Omega ]{}{\operatorname{ess}\sup} \,   u^{+} + & \stackbin[ U_{1,j} \cap \Gamma ]{}{\operatorname{ess}\sup} \, u^{+} \\[3pt]
 \leq & \displaystyle \,  C\left(1+ \max \left\{   \widehat{\fc}_2^{\frac{\qxp _2 \xv }{\qxp _2 (\xv -2) - \xv}}   , \widehat{\fc}_3 ^{\frac{\qxp _3 \xw }{\qxp _3 (\xw -2) - \xw}} \right\}  \right)\\[3pt]
& \displaystyle \, \cdot \left(\Vert u^+ \Vert _{2 , \hat U _{1,j} \cap \Omega   , \wghtv _2}  +   \Vert  u^+ \Vert _{2 , \hat U _{1,j} \cap \Gamma    , \wghtw }  \right) \\[3pt]
& \displaystyle \, + C\left( \Vert f \Vert _{ q_2 ,  \hat U _{1,j} \cap \Omega    , \wghtv _2}  + \Vert f_1 \Vert _{q_3 , \hat U _{1,j} \cap \Gamma  , \wghtw} \right)(1+\widehat{\fc}_2 + \widehat{\fc}_3)^{-1/2},
\end{align*}
where  $C=C (  n , \mathcal C_{\xv} , \mathcal C_{\xw} ,  \xv , \xw , \qxp _2 , \qxp _3 , \ck _2) >0$.
\end{proposition}


\begin{proof}  

We carry out the iteration.  Define for $i\in \mathbb{N} \cup \{0\}$,
$$
 h_{i}:=h_0 + h \left(1-\frac{1}{2^{i}}\right) \leq  h_0 +h \quad \text { and } \quad  R_{i}:= r _{1,j} +\frac{1}{2^{i}} (\hat r _{1,j} -  r _{1,j}).
$$
where  
$$
h_0 = \cen _2 \max \left\{ \widehat{\fc}_2 ^{\frac{\qxp _2 \xv }{\qxp _2 (\xv -2) - \xv}}   , \widehat{\fc}_3 ^{\frac{\qxp _3 \xw }{\qxp _3 (\xw -2) - \xw}} , 1\right\} \max \left\{ \Vert u^+ \Vert _{2 , \hat U _{1,j} \cap \Omega  , \wghtv _2}  ,   \Vert  u^+ \Vert _{2 , \hat U _{1,j} \cap \Gamma , \wghtw }   \right\},
$$
$\cen _2$ is given by Lemma \ref{alcontb0}, and $h>0$ will be determined later.

We have
$$
h_{i}-h_{i-1}=\frac{h}{2^{i}}\quad \text { and } \quad R_{i-1}-R_{i}=\frac{\hat r _{1,j} -  r _{1,j}}{2^{i}} .
$$

Hence, from \eqref{48}, 
\begin{align}
&\Psi  (h_{i}, R _i) \nonumber \\[3pt]
&\quad \ \leq \cen _3 \left[    2^i  + 2^i\frac{\Vert f \Vert _{ q_2 ,  \hat U _{1,j} \cap \Omega  , \wghtv _2}  + \Vert f_1 \Vert _{q_3 , \hat U _{1,j} \cap \Gamma , \wghtw} + (h _0 + h )(1+\widehat{\fc}_2 + \widehat{\fc}_3)^{1/2} }{h}\right]\frac{2^{\epsilon i}}{h^\epsilon} \Psi ^{1 + \epsilon  } (h_{i-1} , R_{i-1}) \nonumber \\[5pt]
&\quad \ \leq 2\cen _3 \left[ \Vert f \Vert _{ q_2 ,  \hat U _{1,j} \cap \Omega  , \wghtv _2}  + \Vert f_1 \Vert _{q_3 , \hat U _{1,j} \cap \Gamma , \wghtw} + (h _0 + h  )(1+\widehat{\fc}_2 + \widehat{\fc}_3)^{1/2}\right] \frac{2^{(1+\epsilon )i}}{h^{1+\epsilon}} \Psi ^{1 + \epsilon  } (h_{i-1} , R_{i-1}), \label{49}
\end{align}

Next we prove inductively for any $i\in \mathbb{N} \cup \{0\}$,
\begin{equation} \label{51}
\Psi  (h_{i} , R_{i}) \leq \frac{\Psi  (h_{0} , R_{0}) }{\gamma^{i}} \quad \text { for some } \quad \gamma>1,
\end{equation}
if $h$ is sufficiently large. It is true for $i=0$. Suppose it is true for $i-1$. We have
\begin{equation} \label{50}
\begin{aligned}
\Psi ^{1 + \epsilon  } (h_{i-1} , R_{i-1}) \leq & \, \displaystyle \left(  \frac{\Psi (h_0, R_0 )}{\gamma^{i-1}}\right) ^{1+\epsilon }\\[3pt]
=  & \, \displaystyle \frac{\Psi ^{\epsilon } (h_0, R_0)}{\gamma^{i \epsilon -(1+\epsilon )}} \cdot  \frac{\Psi (h_0, R_0)}{\gamma^{i}} .
\end{aligned}
\end{equation}

Then, by \eqref{49} and \eqref{50}, we obtain
\begin{align*}
\Psi (h_{i},  R_i) 
 \leq & \, \displaystyle  2\cen _3 \gamma^{1+\epsilon } \frac{  \Vert f \Vert _{ q_2 ,  \hat U _{1,j} \cap \Omega  , \wghtv _2}  + \Vert f_1 \Vert _{q_3 , \hat U _{1,j} \cap \Gamma , \wghtw} +( h _0 + h  )(1+\widehat{\fc}_2 + \widehat{\fc}_3)^{1/2} }{h^{1+\epsilon }} \\[3pt]
 & \, \displaystyle \cdot \Psi ^\epsilon (h_0, R_0)  \frac{2^{i (1+\epsilon )}}{\gamma^{i \epsilon }} \cdot \frac{\Psi ( h_0, R_0 )}{\gamma^{i}} .
\end{align*}

Choose $\gamma$ first such that $\gamma^{\epsilon }=2^{1+\epsilon }$. Note $\gamma>1$. Next, we need
$$
2\cen _3 \gamma^{1+\epsilon } \left(\frac{\Psi (h_0, R_0)}{h}\right)^{\epsilon }   \frac{  \Vert f \Vert _{ q_2 ,  \hat U _{1,j} \cap \Omega  , \wghtv _2}  + \Vert f_1 \Vert _{q_3 , \hat U _{1,j} \cap \Gamma , \wghtw} + (h _0 + h )(1+\widehat{\fc}_2 + \widehat{\fc}_3)^{1/2}  }{h} \leq 1 .
$$

Therefore, we choose
$$
h=C \left[   \left(\Vert f \Vert _{ q_2 ,  \hat U _{1,j} \cap \Omega  , \wghtv _2}  + \Vert f_1 \Vert _{q_3 , \hat U _{1,j} \cap \Gamma , \wghtw}\right)(1+\widehat{\fc}_2 + \widehat{\fc}_3)^{-1/2} + h _0   + \Psi ( h_0 , R_0) \right],
$$
for $C= C (  n , \mathcal C_{\xv} , \mathcal C_{\xw} ,  \xv , \xw , \qxp _2 , \qxp _3 , \ck _2) >0$ sufficiently large. This proves \eqref{51}.


Taking \( i \rightarrow \infty \) in \eqref{51}, we conclude
$$
\Psi (h_0 + h , r_{1,j} )=0.
$$

Hence, 
\begin{align*}
 \stackbin[ U_{1,j }\cap \Omega ]{}{\operatorname{ess}\sup} \,  u^{+} + & \stackbin[ U _{1,j} \cap \Gamma   ]{}{\operatorname{ess}\sup} \, u^{+} \\[3pt]
 \leq & \displaystyle \,  2  (C+1)\left[ \left( \Vert f \Vert _{ q_2 ,  \hat U _{1,j} \cap \Omega  , \wghtv _2}  + \Vert f_1 \Vert _{q_3 , \hat U _{1,j} \cap \Gamma , \wghtw} \right)(1+\widehat{\fc}_2 + \widehat{\fc}_3)^{-1/2}+ h _0   + \Psi ( h_0 , R_0) \right]\\[5pt]
 \leq & \displaystyle \,  2 (C+1)\Bigg[\left( \Vert f \Vert _{ q_2 ,  \hat U _{1,j} \cap \Omega  , \wghtv _2}  + \Vert f_1 \Vert _{q_3 , \hat U _{1,j} \cap \Gamma , \wghtw}\right)(1+\widehat{\fc}_2 + \widehat{\fc}_3)^{-1/2} \\[3pt]
& \displaystyle \, + \cen _2 \max \left\{ \widehat{\fc}_2 ^{\frac{\qxp _2 \xv }{\qxp _2 (\xv -2) - \xv}}   , \widehat{\fc}_3^{\frac{\qxp _3 \xw }{\qxp _3 (\xw -2) - \xw}} , 1 \right\}  \\[3pt]
& \displaystyle \, \cdot \max \left\{ \Vert u^+ \Vert _{2 , \hat U _{1,j} \cap \Omega  , \wghtv _2 }  ,   \Vert  u^+ \Vert _{2 , \hat U _{1,j} \cap \Gamma , \wghtw }   \right\}   \\[3pt]
& \displaystyle \, + \Vert u^+ \Vert _{2 , \hat U _{1,j} \cap \Omega  , \wghtv _2}  +   \Vert  u^+ \Vert _{2 , \hat U _{1,j} \cap \Gamma , \wghtw }   \Bigg].
\end{align*}
This finishes the proof. \end{proof}

Similarly, we have the following result.
\begin{proposition} \label{86}
Suppose $\kp _{1 , j} ^{2 } (\pvr) \wghtv _1 (\pvr )  \leq \ck _{2} \wghtv _2 (\pvr )$ for   a.e. $\pvr  \in \hat U _{1,j}$.    Assume that   $\fc _2, f \in L^{q_2} (\hat U _{1,j} \cap \Omega  ; \wghtv _2)$  for some  $q _2 >  \xv / (\xv -2)  $. Suppose  
\begin{equation*}
\begin{aligned}
 \ops [u,v_{h} ] + \int  _{\Omega } \fc _2 u v_{h}   \,  \dwv _{2}   
 \leq \int  _{\Omega }    f v_{ h } \, \dwv _{2}  ,  
\end{aligned}
\end{equation*}
for all $h\geq 0 $,   $v _{ h }  = u_{ h }  \zeta ^2 $ and $\zeta \in C^{\infty} _0 ( \hat U_{1,j} )$. Then
\begin{align*}
& \stackbin[ U _{1,j} \cap \Omega ]{}{\operatorname{ess}\sup}    u^{+}   \leq  C\left( 1 +  \widehat{\fc}_2 ^{\frac{\qxp _2 \xv }{\qxp _2 (\xv -2) - \xv}}      \right)\Vert u^+ \Vert _{2 , \hat U _{1,j} \cap \Omega   , \wghtv _2}  +  C  \Vert f \Vert _{ q_2 ,  \hat U _{1,j}  \cap \Omega , \wghtv _2}  (1+\widehat{\fc}_2 )^{-1/2} ,
\end{align*}
where  $C=C (  n , \mathcal C_{\xv} , \xv , \qxp _2 , \ck _2  ) >0$.
\end{proposition}




\begin{proposition} \label{53}
 Suppose $\kp _{0 , i} ^{2 } (\pvr) \wghtv _1 (\pvr )  \leq \ck _{2} \wghtv _2 (\pvr )$ for   a.e. $\pvr  \in \hat U _{0,i}$.    Assume that   $\fc _2, f \in L^{q_2} ( \hat U_{0,i}    ; \wghtv _2)$,   for some  $q _2 >  \xv / (\xv -2) $. Suppose  
\begin{equation*}
\begin{aligned}
 \ops [u,v_{h} ] + \int  _{\Omega } \fc _2 u v_{h}  \,  \dwv _{2}    
 \leq \int  _{\Omega}    f v_{ h } \, \dwv _{2}  ,
\end{aligned}
\end{equation*}
for all $h\geq 0 $,   $v _{ h }  = u_{ h }  \zeta ^2 $ and $\zeta \in C^{\infty} _0 ( \hat U_{0,i} )$. Then
\begin{align*}
 &\stackbin[ U_{0,i}]{}{\operatorname{ess}\sup}  u^{+} 
   \leq  C \left( 1+ \|\fc _2 \|_{\qxp _2 , \hat U_{0,i} , \wghtv _2} ^{\frac{\qxp _2 \xv }{\qxp _2 (\xv -2) - \xv}} \right) \Vert u^+ \Vert _{2 , \hat U _{0,i} , \wghtv _2}  +  C  \Vert f \Vert _{ q_2 ,  \hat U _{0,i} , \wghtv _2}   \left(1+\|\fc _2 \|_{\qxp _2 , \hat U_{0,i} , \wghtv _2}  \right)^{-1/2},
\end{align*}
where  $C=C (  n , \mathcal C_{\xv} ,    \xv ,  \qxp _2  , \ck _2)  >0$.
\end{proposition}
 
 %%%%%%%%%%%%%%%%%%%%%%%%%%%%%%%%%%%%%%%%%%%%%%%%%%%%%%%%%%%%%%%%%%%%%%%%%%%%%%%%%%%%%%%%%%%%%%%%%%%%%%%%%%%%%%%%%%%%%%%%%%%%%%%%%%%%%%%%%%%%%%%%%%%%%%%%%%%%%%%%%%%%%%%%%%%%%%%%%%%%%%%%%%%%%%%%%%%%%%%%%%%%%%%%%%%%%%%%%%%%%%%%%%%%%%%%%%%%%%%%%%%%%%%%%%%%%%%%%%%%%%%%%%%%%%%%%%%%%%%%%%%%%%%%%%%%%%%%%%%%%%%%%%%%%%%%%%%%%%%%%%%%%%%%%%%%%%%%%%%%%%%%%%%%%%%%%%%%%%%%%%%%%%%%%%%%%%%%%%%%%%%%%%%%%%%%%%%%%%%%%%%%%
 
\subsection{Proof of Theorem \ref{59}}

{\bf Step 1.} {For every $u\in\womd$, the functional
$$
v\longmapsto\int_\Omega\iwe uv\,\dwv_2
$$
is bounded on $\womd$. Indeed, by H\"older's inequality and \eqref{62},
$$
\left|\int_\Omega\iwe uv\,\dwv_2\right|
\leq
\|\iwe\|_{\xv/(\xv-2),\Omega,\wghtv_2}
\|u\|_{\xv,\Omega,\wghtv_2}
\|v\|_{\xv,\Omega,\wghtv_2}
\leq
\mathcal C_{\xv}^2
\|\iwe\|_{\xv/(\xv-2),\Omega,\wghtv_2}
\|u\|_{1,2,\Omega,\wghtv_0,\wghtv_1}
\|v\|_{1,2,\Omega,\wghtv_0,\wghtv_1}.
$$
Therefore, by the Lax--Milgram theorem, there exists a unique $Au\in\womd$ such that
$$
\ops[Au,v]
=
( \iwe u , v)_{\wghtv_2} 
\quad\forall v\in\womd.
$$
where 
$$
(w , z)_{\wghtv_2} := \int_\Omega wz\,\dwv_2.
$$

The operator $A:\womd\rightarrow\womd$ is bounded and symmetric with respect to the inner product $\ops[\cdot,\cdot]$. Indeed,
$$
\ops[Au,v]=(\iwe u,v)_{\wghtv_2}=(\iwe v,u)_{\wghtv_2}=\ops[u,Av]
\quad\forall u,v\in\womd.
$$
Since $A$ is bounded and defined on all of $\womd$, it is self-adjoint.}

\medskip

{\bf Step 2.} {\it $A$ is compact and has infinitely many positive and negative eigenvalues.}

\medskip

Suppose first that $\spt\iwe\subset\Omega_m$ for some $m\in\mathbb N$. Let $(u_i)$ be bounded in $\womd$. Then
$$
\begin{aligned}
\|Au_i-Au_j&\|_{1,2,\Omega,\wghtv_0,\wghtv_1}^2\\[3pt]
=&\,\ops[A(u_i-u_j),A(u_i-u_j)]\\[5pt]
\leq&\,\|\iwe\|_{L^\infty(\Omega_m)}
\|u_i-u_j\|_{\xv/(\xv-1),\Omega_m,\wghtv_2}
\|Au_i-Au_j\|_{\xv,\Omega_m,\wghtv_2}.
\end{aligned}
$$
From \eqref{62},
$$
\|Au_i-Au_j\|_{1,2,\Omega,\wghtv_0,\wghtv_1}
\leq
\mathcal C_{\xv}\|\iwe\|_{L^\infty(\Omega)}
\|u_i-u_j\|_{\xv/(\xv-1),\Omega_m,\wghtv_2}.
$$
Since the weights are bounded above and below by positive constants on compact subsets of $M$ and
$$
\frac{\xv}{\xv-1}<2,
$$
the embedding
$$
W^{1,2}(\Omega_m;\wghtv_0,\wghtv_1)
\longrightarrow
L^{\xv/(\xv-1)}(\Omega_m;\wghtv_2)
$$
is compact. Hence, $A$ is compact in this case.

In general, set
$$
\iwe_m(\pvr):=\left\{\begin{aligned}
&\iwe(\pvr) &&\text{if }\pvr\in\Omega_m,\\
&0 &&\text{if }\pvr\in\Omega\setminus\Omega_m,
\end{aligned}\right.
$$
and define $A_m:\womd\rightarrow\womd$ by
$$
(\iwe_m u,v)_{\wghtv_2}=\ops[A_mu,v]
\quad\forall u,v\in\womd.
$$
Then
$$
\begin{aligned}
\|A_mu-Au&\|_{1,2,\Omega,\wghtv_0,\wghtv_1}^2\\[3pt]
\leq&\,
\|\iwe_m-\iwe\|_{\xv/(\xv-2),\Omega,\wghtv_2}
\|u\|_{\xv,\Omega,\wghtv_2}
\|A_mu-Au\|_{\xv,\Omega,\wghtv_2}.
\end{aligned}
$$
From \eqref{62},
$$
\|A_mu-Au\|_{1,2,\Omega,\wghtv_0,\wghtv_1}
\leq
C\|\iwe_m-\iwe\|_{\xv/(\xv-2),\Omega,\wghtv_2}
\|u\|_{\xv,\Omega,\wghtv_2}.
$$
Thus, $A_m\rightarrow A$ in operator norm. Since every $A_m$ is compact, $A$ is compact.

%%
%%
%%

We now prove that $A$ has infinitely many eigenvalues of each sign. Choose nonempty open sets
$$
O_+\Subset\{\pvr\in\Omega\mid\iwe(\pvr)>0\},
\qquad
O_-\Subset\{\pvr\in\Omega\mid\iwe(\pvr)<0\}.
$$

%%
%%
%%
%%

For every $k\in\mathbb N$, choose linearly independent functions
$$
\zeta _1^+,\ldots,\zeta _k^+\in C_0^\infty(O_+).
$$
For every nonzero $\zeta \in\operatorname{span}\{\zeta _1^+,\ldots,\zeta _k^+\}$,
$$
\ops[A\zeta ,\zeta ]
=
\int_\Omega\iwe\zeta ^2\,\dwv_2>0.
$$



%%
%%
%%
%%


Therefore, by \cite[Theorem 4.15, Section 4.5]{zbMATH06377290},
$$
\dim\operatorname{Ran}P_A((0,\infty))\geq k.
$$
Since $k$ is arbitrary, the range of the positive spectral
projection $P_A((0,\infty))$ is infinite-dimensional. The same argument in $O_-$ gives
$$
\dim\operatorname{Ran}P_A((-\infty,0))=\infty.
$$
By \cite[Theorem 6.6, Section 6.2]{zbMATH06377290}, the nonzero spectrum of the compact self-adjoint operator $A$ consists of eigenvalues of finite multiplicity whose only possible accumulation point is $0$.


Hence, \eqref{61} has infinitely many eigenvalues
$$
\cdots\leq\lambda_2^-\leq\lambda_1^-<0<\lambda_1^+\leq\lambda_2^+\leq\cdots,
$$
with $\lambda_k^-\rightarrow-\infty$ and $\lambda_k^+\rightarrow+\infty$.

Moreover, as a consequence of Propositions \ref{86} and \ref{53}, provided that the first inequality in \ref{47} is also satisfied, every eigenfunction belongs to \(L^\infty_{\loc}(\overline\Omega)\).

\medskip

{\bf Step 3.} {\it If \ref{47} holds, then $u\in L^\infty(\Omega)$. Moreover, if $u\in W^{1,2}_0(\overline\Omega;\wghtv_0,\wghtv_1)$, then}
\begin{gather}\label{88}
\lim_{m\rightarrow\infty}
\stackbin[\Omega^m]{}{\operatorname{ess}\sup}\,|u|=0.
\end{gather}

Let $\qxp_2>\xv/(\xv-2)$. Applying Propositions \ref{86} and \ref{53} to $u$ and $-u$, with $\fc_2=-\lambda\iwe$ and $f=0$, there exists
$$
C_1=C_1(n,\mathcal C_{\xv},\xv,\qxp_2,\ck_2)>0
$$
such that
\begin{gather*}
\stackbin[U_{1,j}\cap\Omega]{}{\operatorname{ess}\sup}\,|u|
\leq
C_1\left(
1+\|\lambda\iwe\|_{\qxp_2,\hat U_{1,j}\cap\Omega,\wghtv_2}^{\frac{\qxp_2\xv}{\qxp_2(\xv-2)-\xv}}
\right)
\|u\|_{2,\hat U_{1,j}\cap\Omega,\wghtv_2},\\[5pt]
\stackbin[U_{0,i}]{}{\operatorname{ess}\sup}\,|u|
\leq
C_1\left(
1+\|\lambda\iwe\|_{\qxp_2,\hat U_{0,i},\wghtv_2}^{\frac{\qxp_2\xv}{\qxp_2(\xv-2)-\xv}}
\right)
\|u\|_{2,\hat U_{0,i},\wghtv_2}
\end{gather*}
for all $i,j\in\mathbb N$.

From \ref{47},
$$
\|\lambda\iwe\|_{\qxp_2,\hat U_{k,i}\cap\Omega,\wghtv_2}
\leq
|\lambda|\|\iwe\|_{L^\infty(\Omega)}\ck_3^{1/\qxp_2}
$$
and
$$
\|u\|_{2,\hat U_{k,i}\cap\Omega,\wghtv_2}
\leq
\ck_3^{\frac{\xv-2}{2(\xv)}}
\|u\|_{\xv,\hat U_{k,i}\cap\Omega,\wghtv_2}
\leq
\ck_3^{\frac{\xv-2}{2(\xv)}}\mathcal C_{\xv}
\|u\|_{1,2,\Omega,\wghtv_0,\wghtv_1}.
$$
The preceding estimates are uniform in $i$ and $j$. Since the sets $U_{k,i}$ cover $\overline\Omega$, we obtain $u\in L^\infty(\Omega)$.


Let $\epsilon>0$. Since $u\in W^{1,2}_0(\overline\Omega;\wghtv_0,\wghtv_1)$, there exists $\zeta\in C_0^\infty(M)$ such that
$$
\|u-\zeta\|_{1,2,\Omega,\wghtv_0,\wghtv_1}<\epsilon.
$$
Let $m_0\in\mathbb N$ be such that $\spt\zeta\subset D_{m_0}$.

By the local finiteness of $\{\hat U_{k,i}\}$, only finitely many sets $\hat U_{k,i}$ meet $\overline{D_{m_0}}$. Since every corresponding $U_{k,i}$ is relatively compact, there exists $m_1\geq m_0$ such that, whenever $U_{k,i}\cap D^{m_1}\neq\varnothing$, one has
$$
\hat U_{k,i}\cap\spt\zeta=\varnothing.
$$
For these indices, $\zeta=0$ in $\hat U_{k,i}$ and, from \eqref{62},
$$
\|u\|_{\xv,\hat U_{k,i}\cap\Omega,\wghtv_2}
=
\|u-\zeta\|_{\xv,\hat U_{k,i}\cap\Omega,\wghtv_2}
\leq
\mathcal C_{\xv}\epsilon.
$$

\begin{align*}
 \stackbin[ U _{1,j} \cap \Omega ]{}{\operatorname{ess}\sup}  \,  |u| 
   \leq  & \displaystyle \, C_1 \left[  1+  \left(\|\lambda \iwe\|_{L^\infty (\Omega)} ^{\qxp _2} \ck _3 \right) ^{ \frac{\xv}{\qxp _2 (\xv -2) - \xv}} \right] \ck _3 ^{\frac{\xv -2}{\xv}}\Vert u \Vert _{\xv , \hat U_{1,j} \cap  \Omega   , \wghtv _2}  \\[3pt]
 \leq & \displaystyle \,  C_1 \left[ 1+  \left(\|\lambda \iwe\|_{L^\infty (\Omega)} ^{\qxp _2} \ck _3 \right) ^{ \frac{\xv }{\qxp _2 (\xv -2) - \xv}}    \right] \\[3pt]
& \displaystyle \, \cdot  \ck _3 ^{\frac{\xv -2}{\xv}} \left( \Vert u - \zeta \Vert _{\xv ,  \hat U_ {1,j}  \cap \Omega   , \wghtv _2} + \Vert  \zeta \Vert _{\xv ,  \hat U_ {1,j}  \cap \Omega   , \wghtv _2} \right)    \\[5pt]
  \leq & \displaystyle \,  C_1 \left[ 1+  \left(\|\lambda \iwe\|_{L^\infty (\Omega)} ^{\qxp _2} \ck _3 \right) ^{ \frac{\xv }{\qxp _2 (\xv -2) - \xv}}   \right] \mathcal{C} _{\xv}\ck _3 ^{\frac{\xv -2}{\xv}}  \epsilon 
\end{align*}
and 
\begin{align*}
&\stackbin[ U_{0,i}]{}{\operatorname{ess}\sup} \, |u|   \leq  C_1 \left[1+ \left(\|\lambda\iwe\|_{L^\infty (\Omega)} ^{\qxp _2} \ck _3 \right) ^{ \frac{\xv }{\qxp _2 (\xv -2) - \xv}}   \right]\mathcal{C} _{\xv}\ck _3 ^{\frac{\xv -2}{\xv}}       \epsilon.
\end{align*}


Since the right-hand side is independent of $(k,i)$ and the sets $U_{k,i}$ cover $\overline\Omega$, we obtain
$$
\stackbin[\Omega^m]{}{\operatorname{ess}\sup}\,|u|
\leq
C\epsilon
\quad\text{for every }m\geq m_1.
$$
As $\epsilon>0$ is arbitrary, \eqref{88} follows.

This concludes the proof of Theorem \ref{59}.

%%%%%%%%%%%%%%%%%%%%%%%%%%%%%%%%%%%%%%%%%%%%%%%%%%%%%%%%%%%%%%%%%%%%%%%%%%%%%%%%%%%%%%%%%%%%%%%%%%%%%%%%%%%%%%%%%%%%%%%%%%%%%%%%%%%%%%%%%%%%%%%%%%%%%%%%%%%%
%%%%%%%%%%%%%%%%%%%%%%%%%%%%%%%%%%%%%%%%%%%%%%%%%%%%%%%%%%%%%%%%%%%%%%%%%%%%%%%%%%%%%%%%%%%%%%%%%%%%%%%%%%%%%%%%%%%%%%%%%%%%%%%%%%%%%%%%%%%%%%%%%%%%%%%%%%%%

\appendix

\section{Proof of Lemma \ref{alcontb0}}
 Similarly to Lemma \ref{6}, we have the following result.
\begin{lemma} \label{44}
Suppose $0<r<R\leq \hat r _{1,j}$ and fix $j\in \mathbb{N}$. There exists a smooth function $\zeta  : \hat U  _{1,j} \rightarrow [0,1]$ such that  
$$
\zeta  = 1 \text {  in } \psi _{1,j} (B(0,r)),   \quad \spt \zeta  \subset \psi _{1,j} (B(0,R)),
$$
 and 
 $$
 |\nabla \zeta  | _g \leq \cen _1(n)\bnorm d (\psi _{1,j}^{-1}) \bnorm  (R-r)^{-1} \quad \text { in } \quad \hat  U _{1,j}.
 $$ 
\end{lemma}
 
We follow the proof in \cite[Theorem 4.1]{hanlin2011ellipticpartial}; see also \cite[Lemma A.2]{apaza2024yamabecorner}.
\begin{cl} \label{accontb1}
We have
\begin{equation*}
\begin{aligned}
& \int   _{\hat U _{1,j} \cap \Omega }    |\nabla (u_{h_{2}} \zeta  )|^2 \, \dwv _{1}  \\[3pt]
& \quad \ \leq C_1( n) \left[ \frac{(\hat r _{1,j} - r_{1,j} )^2 } {(R-r) ^{2}} \ck _2\int _{\mathscr{U}(h_{2},R)} u_{ h_{2} }^2 \, \dwv _{2}+ \int  _{\hat U _{1,j} \cap \Omega }    g (\nabla u , \nabla v_{ h_{2} }) \, \dwv _{1}  \right].
\end{aligned}
\end{equation*}
\end{cl} 
%\renewcommand{\qedsymbol}{$\square$}
\begin{proof}
We have 
\begin{align*}
\int _{\hat U _{1,j} \cap \Omega } g(\nabla u , \nabla v_{h_{2}}) \, \dwv _{1} = & \displaystyle \, \int _{\hat U _{1,j} \cap \Omega } \left( \zeta  ^2 g(\nabla u , \nabla u_{h_{2}} ) + 2u_{ h_{2} } \zeta  g(\nabla u , \nabla \zeta   )  \right) \, \dwv _{1} \\[3pt]
 \geq & \displaystyle \,   \int _{\hat U _{1,j} \cap \Omega } \zeta  ^2  | \nabla u_{h_{2}} |^2 \, \dwv _{1}   - \int _{\hat U _{1,j} \cap \Omega } 2 \zeta   |\nabla u _{h_{2}}  | \cdot |\nabla \zeta   |  \cdot |u_{h_{2}} |  \, \dwv _{1} \\[5pt]
 \geq & \displaystyle \,  \frac{1}{2}\int _{\hat U _{1,j} \cap \Omega } \zeta  ^2  | \nabla u_{h_{2}} |^2 \, \dwv _{1}   - 2\int _{\hat U _{1,j} \cap \Omega }   |\nabla \zeta   |^2  u_{h_{2}}  ^2 \, \dwv _{1} .
\end{align*}
From Lemma \ref{44} and condition \ref{47}, we conclude the claim.\end{proof}


The inequality \eqref{a168} and  Claim \ref{accontb1} imply 
\begin{equation} \label{a169}
\begin{aligned}
\min \{1, & C_1\}  \ops [u_{h_{2}}  \zeta , u_{h_{2}}  \zeta] \\[3pt]
 \leq & \displaystyle \,  \int  _ {\hat U _{1,j} \cap \Omega }  |\nabla (u_{h_{2}}  \zeta  )|^2 \, \dwv _{1}   + C_1 \int  _{\hat U _{1,j} \cap \Omega }  u v_{h_{2}} \,  \dwv _{0}  \\[3pt]
 \leq  & \displaystyle \,    C_1 \left[ \frac{ (\hat r _{1,j} - r_{1,j})^2}{(R-r) ^{2}} \ck _{2} \int _{\mathscr{U}(h_{2},R)} u_{h_{2}}  ^2 \, { \dwv _{2}} \right. - \int  _{\hat U _{1,j} \cap \Omega } \fc  _2 u v_{h_{2}}   \, \dwv _{2}    \\[5pt]
& \displaystyle \,  \left.      - \int  _{ \hat U _{1,j} \cap \Gamma} \fc _3 u v_{h_{2}}  \, \dww   + \int  _{\hat U _{1,j} \cap \Omega }    f v_{h_{2}}  \, \dwv _{2} + \int_{\hat U _{1,j} \cap \Gamma }  f_1 v _{h_{2}} \, \dww  \right]  .
\end{aligned} 
\end{equation}


Next, we will estimate the terms on the right-hand side of \eqref{a169}.

\begin{cl}\label{accontb2}
The following estimates are valid:

\medskip

\noindent $(i)$ One has
\begin{align*}
-\int _{\hat U _{1,j} \cap \Gamma }  \fc _3  u & (u_{h_{2}}  \zeta ^2 ) \,  \dww    \\[3pt] 
 \leq & \displaystyle \,  2 \mathcal C ^2 _{\xw} \Vert \fc _3 \Vert _{q_3 ,\hat U _{1,j} \cap \Gamma , \wghtw}  \wghtw  ( \{u_{h_{2}}  \zeta \neq 0\} \cap \hat U _{1,j} \cap \Gamma ) ^ {\frac{\xw - 2}{\xw} - \frac{1}{q_3}}  \ops [ u_{h_{2}}  \zeta , u_{h_{2}}  \zeta ] \\[5pt]
& \displaystyle \, +  h_{2}^2 \Vert \fc _3 \Vert _{q_3 , \hat U _{1,j} \cap \Gamma , \wghtw } \wghtw  (\{ u_{h_{2}}  \zeta \neq 0\} \cap \hat U _{1,j} \cap \Gamma )^{1-\frac{1}{q_3}}. 
\end{align*}

\noindent  $(ii)$ For all  $\delta >0$,
\begin{align*}
&\int _{\hat U _{1,j} \cap \Gamma }  f_1 u_{h_{2}}  \zeta ^2  \dww\\[3pt]  
& \quad \ \leq \mathcal C_{\xw} \left( \delta ^{-1}  \wghtw  (\{ u_{h_{2}}  \zeta \neq 0\}\cap \hat U _{1,j} \cap \Gamma ) ^ {  \frac{2(\xw -1)}{\xw} - \frac{2}{q_3}} \Vert f_1 \Vert _{ q_3 , \hat U _{1,j} \cap \Gamma , \wghtw } ^2  + \delta \ops [ u_{h_{2}}  \zeta  , u_{h_{2}}  \zeta ] \right).
\end{align*} 

\medskip

\noindent  $(iii)$ One has
\begin{align*}
-\int _{\hat U _{1,j} \cap \Omega }  \fc _2 u & ( u_{h_{2}}  \zeta ^2 ) \, \dwv _{2} \\[3pt]
 \leq & \displaystyle \,  2 \mathcal C_{\xv} ^2\Vert \fc _2 \Vert _{q_2 , \hat U _{1,j} \cap \Omega  , \wghtv _2}  \wghtv _2( \{u_{h_{2}}  \zeta \neq 0\} ) ^ {\frac{\xv -2}{\xv} - \frac{1}{q_2}} \ops [ u_{h_{2}}  \zeta , u_{h_{2}}  \zeta] \\[5pt]
& \displaystyle \, +  h_{2}^2 \Vert \fc _2 \Vert _{q_2 ,  \hat U _{1,j} \cap \Omega   , \wghtv _2} \wghtv _2 ( \{u_{h_{2}}  \zeta \neq 0\} )^{1-\frac{1}{q_2}}.
\end{align*}

\noindent  $(iv)$ For all $\delta >0$,
\begin{align*}
\int _{\hat U _{1,j} \cap \Omega }  f u_{h_{2}}  & \zeta ^2 \, \dwv _{2} \\[3pt]
 \leq  & \displaystyle \, \mathcal C_{\xv} \left( \delta ^{-1} \wghtv _2 ( \{u_{h_{2}}  \zeta \neq 0\}) ^ {\frac{2(\xv -1)}{\xv} - \frac{2}{q _2}} \Vert f \Vert _{q_2 ,  \hat U _{1,j} \cap \Omega   , \wghtv _2 } ^2  + \delta  \ops [u_{h_{2}}  \zeta , u_{h_{2}}  \zeta ] \right).
\end{align*}

\end{cl}
%\renewcommand{\qedsymbol}{$\square$}
\begin{proof} We start with the following inequalities, which will be used later. Hölder's inequality and \eqref{63} imply
\begin{equation} \label{aecontb1}
\begin{aligned}
 \int _{\hat U _{1,j} \cap \Gamma }  (u_{h_{2}}  \zeta )^2 \, \dww \leq & \displaystyle \,  \wghtw  (\{ u_{h_{2}}  \zeta \neq 0 \} \cap \hat U _{1,j} \cap \Gamma )^{1-\frac{2}{\xw}}\left( \int _{\hat U _{1,j} \cap \Gamma }  | u_{h_{2}}  \zeta |^{\xw} \dww \right)^{\frac{2}{\xw}} \\[3pt]
  \leq  & \displaystyle \,  \mathcal{C} ^2 _{\xw} \wghtw  ( \{ u_{h_{2}}  \zeta \neq 0 \}\cap \hat U _{1,j} \cap \Gamma )^{1-\frac{2}{\xw}} \ops [ u_{h_{2}}  \zeta , u_{h_{2}}  \zeta].
 \end{aligned}
\end{equation}
On the other hand, by  \eqref{62},
\begin{equation} \label{aecontb2}
\begin{aligned}
 \int _{\hat U _{1,j} \cap \Omega } (u_{h_{2}}  \zeta )^2 \, \dwv _{2} \leq & \displaystyle \,  \wghtv _{2}  (\{ u_{h_{2}}  \zeta \neq 0 \}  )^{1-\frac{2}{\xv}}\left( \int _{\hat U _{1,j} \cap \Omega } | u_{h_{2}}  \zeta |^{\xv} \, \dwv _{2} \right)^{\frac{2}{\xv}} \\[3pt]
  \leq & \displaystyle \,  \mathcal{C} ^2 _{\xv} \wghtv _{2}  ( \{ u_{h_{2}}  \zeta \neq 0 \} )^{1-\frac{2}{\xv}} \ops [ u_{h_{2}}  \zeta , u_{h_{2}}  \zeta].
 \end{aligned}
\end{equation}

\medskip

\noindent $(i)$ \, Since $2/\xw + 1/q_3 <1$,
\begin{align*}
-\int _{\hat U _{1,j} \cap \Gamma }  \fc _3  u (u_{h_{2}}  \zeta ^2 )  \,  \dww   = & \displaystyle \,  -\int _{\{ u_{h_{2}}  \zeta  \neq 0\} \cap \hat U _{1,j} \cap \Gamma }  \fc _3 \left(  u_{h_{2}}  ^2 + h_{2}  u_{h_{2}}  \right) \zeta ^2   \dww   \\[3pt]
 \leq & \displaystyle \, 2 \int _{\{ u_{h_{2}}  \zeta  \neq 0 \} \cap \hat U _{1,j} \cap \Gamma }  |\fc _3|  u_{h_{2}}  ^2 \zeta ^2  \, \dww   + h_{2}^2    \int _{\{ u_{h_{2}}  \zeta  \neq 0 \} \cap \hat U _{1,j} \cap \Gamma }     |\fc _3|\zeta ^2  \, \dww   \\[3pt]
 \leq & \displaystyle \,  2\left( \int _{\hat U _{1,j} \cap \Gamma }  |\fc _3| ^{q_3}  \dww  \right) ^{\frac{1}{q_3}} \left( \int _{\{ u_{h_{2}}  \zeta \neq 0 \} \cap \hat U _{1,j} \cap \Gamma }   | u_{h_{2}}  \zeta |^ {\xw }  \dww    \right) ^{\frac{2}{\xw}}  \\[3pt]
& \displaystyle \, \cdot \wghtw  ( \{ u_{h_{2}}  \zeta \neq 0 \} \cap \hat U _{1,j} \cap \Gamma  ) ^ {1-\frac{2}{\xw} - \frac{1}{q_3}} \\[3pt]
& \displaystyle \,  + h_{2}^2 \Vert \fc _3 \Vert _{q_3 , \hat U _{1,j} \cap \Gamma , \wghtw}  \wghtw  (  \{ u_{h_{2}}  \zeta \neq 0 \} \cap \hat U _{1,j} \cap \Gamma )^{1-\frac{1}{q_3}}\\[7pt]
 \leq   & \displaystyle \, 2 \mathcal C ^2 _{\xw} \Vert \fc _3 \Vert _{ q_3 , \hat U _{1,j} \cap \Gamma , \wghtw}  \wghtw  (  \{ u_{h_{2}}  \zeta \neq 0 \} \cap \hat U _{1,j} \cap \Gamma ) ^ {\frac{\xw -2}{\xw} - \frac{1}{q_3}} \ops [ u_{h_{2}}  \zeta , u_{h_{2}}  \zeta] \\[3pt]
& \displaystyle \,  + h_{2}^2 \Vert \fc _3 \Vert _{q_3 , \hat U _{1,j} \cap \Gamma , \wghtw  } \wghtw  (  \{ u_{h_{2}}  \zeta \neq 0 \} \cap \hat U _{1,j} \cap \Gamma )^{1-\frac{1}{q_3}}, 
\end{align*}
by \eqref{aecontb1}. This concludes the proof of $(i)$.\\

\noindent $(ii)$ \ Since      $1/q_3 + 1/\xw < 1$  and  $0\leq \zeta \leq 1$,
\begin{align*}
& \int _{\hat U _{1,j} \cap \Gamma }  f_1 u_{h_{2}}  \zeta ^2  \dww    \\[3pt]
& \quad \  \leq   \left( \int _{\hat U _{1,j} \cap \Gamma }  |f_1| ^{q_3}  \dww   \right) ^{\frac{1}{q_3}} \left( \int _{\hat U _{1,j} \cap \Gamma }   | u_{h_{2}}  \zeta |^ {\xw}  \dww    \right) ^{ \frac{1}{\xw} }  \wghtw  (\{ u_{h_{2}}  \zeta \neq 0\}\cap \hat U _{1,j} \cap \Gamma ) ^ {1- \frac{1}{q_3} - \frac{1}{\xw}}.
\end{align*}
By \eqref{aecontb1},
\begin{align*}
&\int _{\hat U _{1,j} \cap \Gamma }  |f_1| u_{h_{2}}  \zeta ^2   \dww    \\[3pt]
& \quad \ \leq  \mathcal C_{\xw} \Vert f_1 \Vert _{q_3 , \hat U _{1,j} \cap \Gamma , \wghtw }  \sqrt{\ops [ u_{h_{2}}  \zeta , u_{h_{2}}  \zeta ]}   \wghtw  (\{ u_{h_{2}}  \zeta \neq 0\}\cap \hat U _{1,j} \cap \Gamma ) ^ {\frac{\xw -1 }{\xw} - \frac{1}{q_3}} \\[5pt]
&\quad \ \leq  \mathcal C_{\xw} \delta ^{-1} \wghtw  ( \{ u_{h_{2}}  \zeta \neq 0\}\cap \hat U _{1,j} \cap \Gamma ) ^ {\frac{2(\xw -1)}{\xw} - \frac{2}{q_3}}  \Vert f_1 \Vert _{q_3 , \hat U _{1,j} \cap \Gamma , \wghtw} ^2  +  \mathcal C_{\xw} \delta \ops [ u_{h_{2}}  \zeta , u_{h_{2}}  \zeta ].
\end{align*}
This proves $(ii)$. The proofs of $(iii)$ and $(iv)$ follow the same lines as in $(i)$ and $(ii)$. \end{proof}

Let us observe that  
\begin{gather*}
1- \frac{ 1}{q_2} \leq \frac{2(\xv -1) }{ \xv} - \frac{2}{q_2}, \quad   1- \frac{1}{q_3} \leq \frac{2(\xw -1) }{ \xw} - \frac{2}{q_3},\\[3pt]  
\{u_{h_{2}}  \zeta  \neq 0\} \subset \mathscr{U}(h_{2},R),  \quad \{u_{h_{2}}  \zeta  \neq 0\} \cap \hat U _{1,j} \cap \Gamma \subset \mathscr{U} _{\Gamma }  (h_{2},R),\\[3pt]
  \wghtv _{2} (\mathscr{U} (h_{2},R) ) \leq h_{2}^{-2} \int _{\mathscr{U} (h_{2},R)} (u^+)^2 \, \dwv _{2} , \quad   \text {  and  } \quad \wghtw ( \mathscr{U} _{\Gamma }  (h_{2},R) ) \leq h_{2}^{-2} \int _{\mathscr{U} _{\Gamma } (h_{2},R)} (u^+)^2  \, \dww  .
  \end{gather*} 

Hence, by \eqref{a169} and Claim \ref{accontb2}, there exists a constant $N = N (  n , \mathcal C_{\xv} , \mathcal C_{\xw} , \xv , \xw , \qxp _2 , \qxp _3 ) >0 $ such that if  
$$
h_{2}\geq N \max \left\{   \|\fc _2 \|_{\qxp _2 , \hat U_{1,j} \cap \Omega , \wghtv _2} ^{\frac{\qxp _2 \xv }{\qxp _2 (\xv -2) - \xv}}   , \|\fc _3 \|_{\qxp _3 , \Gamma \cap \hat U_{1,j}, \wghtw} ^{\frac{\qxp _3 \xw }{\qxp _3 (\xw -2) - \xw}} , 1\right\}  \max \left\{ \Vert u^+ \Vert _{2 , \hat U _{1,j} \cap \Omega  , \wghtv _2 }  ,   \Vert  u^+ \Vert _{2 , \hat U _{1,j} \cap \Gamma , \wghtw }   \right\},
$$
then
\begin{equation}
\wghtv _2(\mathscr{U}(h_{2},R))<1 , \quad   \wghtw  ( \mathscr{U} _{\Gamma }  (h_{2},R)) \ <1 \label{aecontb5}
\end{equation}
and
\begin{equation} \label{aecontb4}
\begin{split} 
C_{2} \ops [ u_{h_{2}}  \zeta , &  u_{h_{2}}  \zeta] \\[3pt]
 \leq  & \displaystyle \,  \frac{ (\hat r _{1,j} - r_{1,j})^2}{(R-r) ^{2}} \ck _{2} \int _{\mathscr{U}(h_{2},R)} u_{h_{2}}  ^2   \, \dwv _{2}  +    \left(  \Vert f \Vert _{q_2 , \hat U _{1,j} \cap \Omega  , \wghtv _2 } ^2 + h_{2}^2 \widehat{\fc}_2 \right)  \wghtv _2 ( \{u_{h_{2}}  \zeta \neq 0\})^{1-\frac{1}{q_2}}  \\[5pt]
 & \displaystyle \, + \left(  \Vert f_1 \Vert _{q_3 ,\hat U _{1,j} \cap \Gamma , \wghtw} ^2 + h_{2}^2 \widehat{\fc}_3 \right) \wghtw (\{ u_{h_{2}}  \zeta \neq 0\}\cap \hat U _{1,j} \cap \Gamma)^{1-\frac{1}{q_3}} ,
\end{split}
\end{equation}
where $C_{2}=C_{2}(  n , \mathcal C_{\xv} , \mathcal C_{\xw}  )>0$, and
$$
\widehat{\fc}_2 :=\Vert \fc _2 \Vert _{q_2 , \hat U _{1,j} \cap \Omega  , \wghtv _2 } ,\qquad \widehat{\fc}_3 :=\Vert \fc _3 \Vert _{q_3 , \hat U _{1,j} \cap \Gamma  , \wghtw } .
$$


From   \eqref{aecontb2} and   \eqref{aecontb4}, we have
\begin{equation}\label{aecontb6}
\begin{split} 
C_{3} \int   _{\hat U _{1,j} \cap \Omega } &    (u_{h_{2}}  \zeta )^2 \, \dwv _{2}   \\[3pt]
 \leq & \displaystyle \,    \frac{ (\hat r _{1,j} - r_{1,j})^2}{(R-r) ^{2}} \wghtv _2 ( \{u_{h_{2}}  \zeta \neq 0\})^{1- \frac{2}{\xv}}  \left( \int _{\mathscr{U}(h_{2},R)} u_{h_{2}}  ^2 \, \dwv _{2}    + \int _{\mathscr{U} _{\Gamma } (h_{2},R)}  u_{h_{2}}  ^2 \, \dww   \right)   \\[3pt]
 & \displaystyle \, + \left[  \Vert f \Vert _{q_2 , \hat U _{1,j} \cap \Omega  , \wghtv _2 } + \Vert f_1 \Vert _{q_3 , \hat U _{1,j} \cap \Gamma , \wghtw } + h_{2} (\widehat{\fc}_2+\widehat{\fc}_3 )^{1/2} \right]^2 \\[3pt]
& \displaystyle \,  \cdot   \left( \wghtv _2 ( \{u_{h_{2}}  \zeta \neq 0\})^{1- \frac{2}{\xv} + 1 -\frac{1}{q_2} }  + \wghtv _2 ( \{u_{h_{2}}  \zeta \neq 0\})^{1- \frac{2}{\xv}}  \wghtw  ( \{ u_{h_{2}}  \zeta \neq 0\}\cap \hat U _{1,j} \cap \Gamma )^{1-\frac{1}{q_3}} \right) .
\end{split}
\end{equation}
and, by \eqref{aecontb1},
\begin{equation}\label{aecontb7}
\begin{split}
C_{4} \int _{\hat U _{1,j} \cap \Gamma } & ( u_{h_{2}}  \zeta)^2   \dww   \\[3pt]
 \leq & \displaystyle \,    \frac{ (\hat r _{1,j} - r_{1,j})^2}{(R-r) ^{2}}    \wghtw  ( \{ u_{h_{2}}  \zeta \neq 0\} \cap \hat U _{1,j} \cap \Gamma )^{1 - \frac{2}{\xw}} \left( \int _{\mathscr{U} (h_{2}, R )} u_{h_{2}}  ^2  \, \dwv _{2}   + \int _{\mathscr{U}_{\Gamma } (h_{2},R)} u_{h_{2}}  ^2  \dww   \right)       \\[3pt]
& \displaystyle \, + \left[  \Vert f \Vert _{q _2 , \hat U _{1,j} \cap \Omega   , \wghtv _2} +  \Vert f_1 \Vert _{q_3 , \hat U _{1,j} \cap \Gamma , \wghtw}  + h_{2} (\widehat{\fc}_2+\widehat{\fc}_3 )^{1/2}\right]^2 \\[3pt]
& \displaystyle \, \cdot  \left(\wghtw  (\{ u_{h_{2}}  \zeta \neq 0\} \cap \hat U _{1,j} \cap \Gamma )^{1 - \frac{2}{\xw}}  \wghtv _{2}  (\{u_{h_{2}}  \zeta \neq 0\} )^{1 -\frac{1}{q_2}}   \right. \\[3pt]
& \displaystyle \, \left. +   \wghtw  ( \{ u_{h_{2}}  \zeta \neq 0\} \cap \hat U _{1,j} \cap \Gamma )^{1 - \frac{2}{\xw}+ 1 -\frac{1}{q_3}} \right) ,
 \end{split}
\end{equation}
where $C_i = C_i (  n , \mathcal C_{\xv} , \mathcal C_{\xw} , \ck _{2} )>0$, $i=3, 4$.

Set
$$
\epsilon  := \min \left\{ 1-\frac{1}{q_2} -\frac{2}{\xw} , 1-\frac{1}{q_3} -\frac{2}{\xv} , 1-\frac{1}{q_2} -\frac{2}{\xv} , 1-\frac{1}{q_3} -\frac{2}{\xw}  \right\}>0
$$


 By Young's inequality,
\begin{equation}\label{a206}
\begin{aligned}
& \wghtv _2 (\{u_{h_{2}}  \zeta \neq 0\} )^{1-\frac{2}{\xv}}  \wghtw (\{ u_{h_{2}}  \zeta \neq 0\}  \cap \hat U _{1,j} \cap \Gamma ) ^{1 -\frac{1}{q_3}}\\[3pt] 
& \quad \ \leq  \frac{1}{C_{5 }}  \wghtv _2 ( \{u_{h_{2}}  \zeta \neq 0\} )^{1+\epsilon } +  \frac{C_{5 } -1}{C_{5 }} \wghtw ( \{ u_{h_{2}}  \zeta \neq 0\} \cap \hat U _{1,j} \cap \Gamma ) ^{\left(1-\frac{1}{q_3}\right)\frac{C_{5 }}{C_{5 }-1} }     ,
\end{aligned}
\end{equation}
\begin{equation} \label{a254}
\begin{aligned}
& \wghtv _2 ( \{u_{h_{2}}  \zeta \neq 0\} )^{1 -\frac{1}{q_2}}  \wghtw ( \{ u_{h_{2}}  \zeta \neq 0\} \cap \hat U _{1,j} \cap \Gamma )^{1- \frac{2}{\xw}}\\[3pt] 
& \quad \leq \frac{C_{6 }-1}{C_{6 }} \wghtv _2 ( \{ u_{h_{2}}  \zeta \neq 0\}  ) ^{\left(1-\frac{1}{q_2}\right)\frac{C_{6 }}{C_{6 }-1} }  + \frac{1}{C_{6 }}  \wghtw ( \{u_{h_{2}}  \zeta \neq 0\} \cap \hat U _{1,j} \cap \Gamma ) ^{1+\epsilon },
\end{aligned}
\end{equation}
where $C_{5 } (1- 2/\xv ) =1+\epsilon $ and $C_{6 }(1-2/\xw)= 1+\epsilon$. Observe also that 
$$
\left(1-\frac{1}{q_3}\right)\frac{C_{5 }}{C_{5 }-1} \geq 1+\epsilon  \ \text { and  } \ \left(1-\frac{1}{q_2}\right)\frac{C_{6 }}{C_{6 }-1}\geq 1+\epsilon .
$$


 
The inequalities \eqref{aecontb5}, \eqref{aecontb6} - \eqref{a254} imply
\begin{equation}\label{a151}
\begin{split}
C_{7} \Psi ^2 (h_{2} , r ) \leq & \, \displaystyle   \frac{ (\hat{r}_{1, j}-r_{1, j})^2}{(R-r)^2} \left(   \wghtv _2 (\{u_{h_{2}}  \zeta \neq 0 \})  +\wghtw ( \{u_{h_{2}}  \zeta \neq 0 \}\cap \hat U _{1,j} \cap \Gamma) \right)^{\epsilon }\Psi ( h_{2}, R ) ^2 \\[3pt]
& \, \displaystyle + \left[  \Vert f \Vert _{q_2 , \hat U _{1,j} \cap \Omega  , \wghtv _2 } +  \Vert f_1 \Vert _{q_3 , \hat U _{1,j} \cap \Gamma , \wghtw}  + h_{2} (\widehat{\fc}_2+\widehat{\fc}_3 )^{1/2} \right]^2 \\[3pt]
& \, \displaystyle \cdot \left( \wghtv _2 (\{u_{h_{2}}  \zeta \neq 0 \})  +\wghtw ( \{u_{h_{2}}  \zeta \neq 0 \}\cap \hat U _{1,j} \cap \Gamma) \right)^{1+\epsilon },
\end{split}
\end{equation}
where $C_{7}= C_{7} (  n , \mathcal C_{\xv} , \mathcal C_{\xw}  , \xv , \xw , \qxp _2 , \qxp _3 , \ck _{2}) >0$. Consider \eqref{a151} and the following claim, which is proved as in \cite{hanlin2011ellipticpartial}:
\begin{cl}
If $h_{2}>h_{1} $, then
\begin{gather*}
\int _{\mathscr{U}(h_{2},R)} u_{h_{2}}  ^2 \,  \dwv _2  \leq \int _{\mathscr{U} (h_{1} ,R)} u_{h_{1}}   ^2 \, \dwv _2  , \quad \int _{\mathscr{U} _{\Gamma } (h_{2},R)}  u_{h_{2}}  ^2   \dww  \leq \int _{\mathscr{U} _{\Gamma }  (h_{1} ,R)} u_{h_{1}}   ^2  \dww  ,\\[3pt]
 \wghtv _2 ( \{u_{h_{2}}  \zeta  \neq 0\} ) \leq  \frac{1}{(h_{2}-h_{1} )^2}\int _{\mathscr{U}(h_{1} ,R)} u_{h_{1}}   ^2\, \dwv _2, \\[3pt]
  \wghtw ( \{ u_{h_{2}}  \zeta \neq 0\} \cap \hat U _{1,j} \cap \Gamma ) \leq  \frac{1}{(h_{2}-h_{1} )^2}\int _{\mathscr{U}_{\Gamma } (h_{1} ,R)} u_{h_{1}}   ^2  \dww .
\end{gather*}
\end{cl}
\noindent Therefore,
\begin{align*}
& C_7 \Psi ^2 (h_{2},r) \\
& \quad \ \leq  \left\{ \frac{ (\hat r _{1,j} - r_{1,j})^2}{(R-r) ^{2}(h_{2}-h_{1} )^{2\epsilon }}  + \frac{  [  \Vert f \Vert _{q_2  ,  \hat U _{1,j} \cap \Omega  , \wghtv _2}  +  \Vert f_1 \Vert _{q_3 , \hat U _{1,j} \cap \Gamma , \wghtw} + h_{2} (\widehat{\fc}_2+\widehat{\fc}_3 )^{1/2}   ]^2}{(h_{2}-h_{1} )^{2(1+\epsilon )}}\right\} \Psi ^{2(1 + \epsilon  )} (h_{1} ,R).
\end{align*}
This proves Lemma \ref{alcontb0}.





%%%%%%%%%%%%%%%%%%%%%%%%%%%%%%%%%%%%%%%%%%%%%%%%%%%%%%%%%%%%%%%%%%%%%%%%%%%%%%%%%%%%%%%%%%%%%%%%%%%%%%%%%%%%%%%%%%%%%%%%%%%%%%%%%%%%%%%%%%%%%%%%%%%%%%%%%%%%%%%%%%%%%%%%%%%%%%%%%%%%%%%%%%%%%%%%%%%%%%%%%%%%%%%%%%%%%%%%%%%%%%%%%%%%%%%%%%%%%%%%%%%%%%%%%%%%%%%%%%%%%%%%%%%%%%%%%%%%%%%%%%%%%%%%%%%%%%%%%%%%%%%%%%%


\vspace{1cm}


%\noindent {\bf Author contributions:} All authors have contributed equally to this work for writing, review and editing. All authors have read and agreed to the published version of the manuscript.

\noindent {\bf Funding:}This work was supported by the Instituto Nacional de Ciência e Tecnologia de Matemática (INCTMat) through CNPq Grant No. 170245/2023-3, and by CNPq-Brazil under Grants Nos. 153232/2024-2 and 150680/2025-2.


\noindent {\bf Data Availability:} No data were used for the research described in the article.

\noindent {\bf Declarations}

\noindent {\bf Conflict of interest:} The author declares no conflict of interest.


%\noindent {\bf Declaration of generative AI and AI-assisted technologies in the writing process:} During the preparation of this work, the authors used the AI ChatGPT and the AI Gemini to correct grammar and orthography in the text. After using this tool/service, the authors reviewed and edited the content as needed and takes full responsibility for the content of the publication.


%%%%%%%%%%%%%%%%%%%%%%%%%%%%%%%%%%%%%%%%%%%%%%%%%%%%%%%%%%%%%%%%%%%%%%%%%%%%%%%%%%%%%%%%%%%%%%%%%%%%%%%%%%%%%%%%%%%%%%%%%%%%%%%%%%%%%%%%%%%%%%%%%%%%%%%%%%%%%%%%%%%%%%%%%%%%%%%%%%%%%%%%%%%%%%%%%%%%%%%%%%%%%%%%%%%%%%%%%%%%%%%%%%%%%%%%%%%%%%%%%%%%%%%%%%%%%%%%%%%%%%%%%%%%%%%%%%%%%%%%%%%%%%%%%%%%%%%%%%%%%%%%%%%%%%%%%%%%%%%%%%%%%%%%%%%%%%%%%%%%%%%%%%%%%%%%%%%%%%%%%%%%%%%%
 


%\bibliographystyle{plain}
 
 \bibliographystyle{abbrv}
 
    
    \bibliography{ref}


\end{document}
\newpage

```latex
\subsection{Context and Related Work}

On a closed Riemannian manifold, the Laplace--Beltrami operator has
discrete spectrum. For an open subset of \(\mathbb R^n\), discreteness
depends on the chosen realization of the Laplacian and on the boundary
condition. In particular, the Neumann Laplacian has discrete spectrum
whenever the embedding
\[
W^{1,2}(\mathcal M)\longrightarrow L^2(\mathcal M)
\]
is compact. In general, the spectrum of \(\Delta_{\mathcal M}\) need
not be discrete. Several nonstandard geometric settings have therefore
been studied. Conditions for discreteness of the spectrum of the
Laplacian on noncompact complete Riemannian manifolds with specific
geometric structures appear in
\cite{baider1979discretespectra,brooks1984fiitevolume,
bruning1989discretespectrum,donnelly1979purepointspectrum,
escobar1986spectrum,kleine1988discretenessconditions}.

{\color{red}
For weighted Sobolev spaces on unbounded Euclidean domains, compactness
cannot be obtained only from local Sobolev inequalities. A bounded
sequence may preserve its norm while its support moves toward infinity.
Thus, continuous and compact embeddings depend on the interaction
between the weights and the geometry of the domain at infinity. Adams
\cite{adams1971compactembedding} and Gurka \& Opic
\cite{gruka1991weighteiii} study this problem through weighted
embedding conditions on unbounded domains. Gol'dshtein \& Ukhlov
\cite{goldshtein2007weighted} obtain sufficient conditions for
continuous and compact embeddings on smooth domains and on domains
with boundary singularities. Their method uses generalized
quasiconformal mappings to reduce the weighted embedding problem to an
embedding problem on a regular reference domain.
}

{\color{red}
The approach used below is different. We do not assume a global
transformation of \(\Omega\) onto a model domain. The assumptions are
imposed directly on a locally finite family of interior and boundary
charts. The chart radii, the metric distortion, the multiplicity of the
covering, and the local behavior of the weights determine the constants
in the Sobolev and trace inequalities. Compactness is obtained by
requiring the corresponding constants to tend to zero along the
noncompact part of \(\Omega\).
}

{\color{red}
Extension and trace results are related, but they are distinct.
Chua \cite{chua1992extensiontheorems} studies extension operators in
weighted Sobolev spaces, whereas Pflüger
\cite{pfluger1998compactraces} studies compact weighted traces.
In the broader setting of Besov and Triebel--Lizorkin spaces on bounded
Lipschitz domains, Shi \& Yao \cite{shiyao2024rychkov} obtain
quantitative estimates for Rychkov's universal extension operator.
In the present work, the trace inequalities are established directly
in boundary charts. The extension result is therefore independent of
the proofs of the weighted embeddings and traces.
}

Cianchi \& Maz'ya \cite{cianchi2011spectrumnoncomriem} prove that, for
a noncompact Riemannian manifold \(\mathcal M^n\), with \(n\geq2\) and
\(\mathcal H^n(\mathcal M)<\infty\), the embedding
\[
W^{1,2}(\mathcal M)\longrightarrow L^2(\mathcal M)
\]
is compact if and only if
\[
\lim_{s\rightarrow0}\frac{s}{\mu_{\mathcal M}(s)}=0,
\]
which is equivalent to discreteness of the spectrum of
\(\Delta_{\mathcal M}\). Here, the isocapacitary function
\[
\mu_{\mathcal M}:
\left[0,\frac{\mathcal H^n(\mathcal M)}{2}\right]
\longrightarrow[0,\infty]
\]
is defined by
\[
\begin{aligned}
\mu_{\mathcal M}(s):=
\inf\Big\{C(E,G)\mid
&E\text{ and }G\text{ are measurable subsets of }\mathcal M,\\
&E\subset G\subset\mathcal M,\qquad
s\leq\mathcal H^n(E),\qquad
\mathcal H^n(G)\leq
\frac12\mathcal H^n(\mathcal M)
\Big\},
\end{aligned}
\]
where
\[
\begin{aligned}
C(E,G):=\inf\Bigg\{
\int_{\mathcal M}|\nabla u|^2\,\dx
\ \Bigg|\ 
&u\in W^{1,2}(\mathcal M),\qquad
u\geq1\text{ in }E,\\
&u\leq0\text{ in }\mathcal M\setminus G
\text{ up to a set of standard capacity zero}
\Bigg\}.
\end{aligned}
\]

Furthermore, Cianchi \& Maz'ya
\cite{cianchi2013boundseigenfunctions} prove that, if
\(\mathcal H^n(\mathcal M)<\infty\) and
\[
\int_0\frac{\textnormal{d}s}{\mu_{\mathcal M}(s)}<\infty,
\]
then, for every eigenvalue \(\gamma\) of \(\Delta_{\mathcal M}\), there
exists a constant \(C=C(\mu_{\mathcal M},\gamma)\) such that
\[
\|u\|_{L^\infty(\mathcal M)}
\leq
C\|u\|_{L^2(\mathcal M)}
\]
for every eigenfunction \(u\) associated with \(\gamma\).

{\color{red}
The conditions of Cianchi \& Maz'ya are intrinsic and global, since
they are formulated through capacity. Our conditions are local and
quantitative. We apply Euclidean inequalities in each chart and then
sum the resulting estimates by using the finite multiplicity of the
covering. The behavior at infinity is controlled by the quantities
\(\mathcal B_{\qxp,\qxp_0}^m\) and
\(\mathcal B_{\Gamma,\qxp,\qxp_1}^m\).
Their convergence to zero prevents the Sobolev mass from escaping
through the noncompact part of the domain.
}

{\color{red}
A related compactness mechanism is studied by Skrzypczak \& Tintarev
\cite{skrzypczak2020compact}. They analyze compact subsets and
subspaces of Sobolev spaces on noncompact Riemannian manifolds by means
of discretizations and local vanishing conditions. Their discussion
includes the compactness produced by domains that become sufficiently
thin at infinity, by weights, and by symmetry restrictions. The
present work treats the loss of compactness at infinity through
weighted chart estimates and also includes boundary traces.
}

Let \(\mathcal M_0\) be a complete Riemannian manifold, and let
\(\mathsf m\geq0\) be a locally integrable potential. Ouhabaz
\cite{ouhabaz2001spectral} studies the positivity of the \(L^2\)
spectral lower bound of the Schrödinger operator
\[
-\Delta_{\mathcal M_0}+\mathsf m.
\]
He proves that, if \(\mathcal M_0\) satisfies local
\(L^2\)-Poincaré inequalities, namely,
\[
\int_{B(\pvr,r)}
\left|u-\overline{u_{B(\pvr,r)}}\right|^2\,\dvg
\leq
C(R)r^2
\int_{B(\pvr,r)}|\nabla u|^2\,\dvg
\]
for every \(u\in C^\infty(B(\pvr,r))\), \(\pvr\in\mathcal M_0\), and
\(0<r\leq R\), and if \(\mathcal M_0\) satisfies a local doubling
property, then
\[
\inf\sigma(-\Delta_{\mathcal M_0}+\mathsf m)>0,
\]
provided that
\[
\inf_{\pvr\in\mathcal M_0}
\frac{1}{|B(\pvr,r)|}
\int_{B(\pvr,r)}\mathsf m\,\dvg
>0
\]
for some \(r>0\). Ouhabaz also shows that this mean condition is
necessary under additional geometric assumptions on \(\mathcal M_0\).

{\color{red}
Berestycki \& Rossi \cite{zbMATH06441399} study
generalized principal eigenvalues for elliptic operators on unbounded
domains. They show that several properties of the principal eigenvalue
on bounded domains may fail in the unbounded case and relate different
notions of generalized principal eigenvalue to the maximum principle
and to the existence of positive eigenfunctions. Their approach does
not assume compactness of the resolvent. The problem studied below is
different: compactness is recovered through weighted Sobolev
embeddings, and the weight in the eigenvalue equation may change sign.
}

If \(\Sigma_0\) is a bounded domain in \(\mathbb R^n\) with smooth
boundary, the eigenvalue problem
\[
\left\{
\begin{aligned}
-\Delta u&=\lambda u
&&\text{in }\Sigma_0,\\
u&=0
&&\text{on }\partial\Sigma_0,
\end{aligned}
\right.
\]
has an infinite sequence of positive eigenvalues
\[
0<\lambda_1<\lambda_2\leq\cdots\leq\lambda_k\leq\cdots,
\qquad
\lambda_k\longrightarrow\infty
\quad\text{as }k\longrightarrow\infty,
\]
each of finite multiplicity. The same properties hold for
\[
\left\{
\begin{aligned}
-\Delta u+a_0(x)u&=\lambda m(x)u
&&\text{in }\Sigma_0,\\
u&=0
&&\text{on }\partial\Sigma_0,
\end{aligned}
\right.
\]
when \(a_0\) and \(m\) are positive and sufficiently regular on
\(\overline{\Sigma_0}\). For a detailed study of principal eigenvalues
for second-order differential operators that are not necessarily in
divergence form, see Fleckinger, Hernández, \& Thélin
\cite{fleckinger2004existence}.

The existence of principal eigenvalues with positive eigenfunctions is
important in nonlinear problems for which positive solutions are
sought. This occurs, for example, in reaction--diffusion systems
arising in population dynamics, chemical reactions, and combustion;
see \cite{smoller2012shock}.

The classical reference for this theory is the book by Courant \&
Hilbert \cite{courant1962methods}. The main tool is the spectral theory
of compact self-adjoint operators on Hilbert spaces. The variational
characterization yields monotonicity with respect to the coefficients
and, under domain inclusion, with respect to the domain. Continuity
with respect to coefficients or domains requires an appropriate notion
of convergence.

The theory also extends to unbounded coefficients and indefinite
weights. Assume that \(a_0>0\) in \(\Sigma_0\),
\(a_0,m\in L^r(\Sigma_0)\) for some \(r>n/2\), and that the positive
part of \(m\) is nontrivial. Then there exists a unique positive
principal eigenvalue \(\lambda_1^+>0\) with a positive eigenfunction.
If the negative part of \(m\) is nontrivial, there is also a unique
negative principal eigenvalue \(\lambda_1^-<0\) with a positive
eigenfunction; see
\cite{figuereido1982semilinearelliptic,brown1980existence,
weinberger1974variational}.

{\color{red}
For indefinite-weight problems on unbounded domains, the change of sign
of the weight and the possible loss of compactness at infinity must be
treated simultaneously. Allegretto
\cite{allegreto1992eigenvalueindefiniteweight} studies principal
eigenvalues for indefinite-weight elliptic problems in
\(\mathbb R^n\). In the present work, the continuous weighted embedding
controls the bilinear form containing the indefinite weight, while the
compact weighted embedding gives a compact self-adjoint solution
operator. The positive and negative parts of the weight then yield
positive and negative eigenvalues.
}

{\color{red}
The decay of the corresponding eigenfunctions is a separate issue.
Compactness gives the spectral decomposition, but it does not by itself
give pointwise decay. Here, the decay at infinity is obtained by
combining local \(L^\infty\)-estimates with the fact that the weighted
Sobolev norm tends to zero in the noncompact part of \(\Omega\).
Thus, the geometry of the charts and the behavior of the weights enter
both the compactness argument and the final decay estimate.
}

{\color{red}
The results below therefore connect three problems: weighted Sobolev
embeddings and traces on unbounded domains, Sobolev spaces on
noncompact Riemannian manifolds, and elliptic eigenvalue problems with
indefinite weights. The assumptions are stated directly in terms of the
charts and the weights, without requiring curvature bounds, a uniform
positive injectivity radius, or an isocapacitary description of the
domain.
}
```

\end{document}


```bibtex
@misc{goldshtein2007weighted,
  author        = {Gol'dshtein, V. and Ukhlov, A.},
  title         = {Weighted {Sobolev} Spaces and Embedding Theorems},
  year          = {2007},
  eprint        = {math/0703725},
  archivePrefix = {arXiv},
  primaryClass  = {math.FA},
  doi           = {10.48550/arXiv.math/0703725}
}

@misc{skrzypczak2020compact,
  author        = {Skrzypczak, Leszek and Tintarev, Cyril},
  title         = {On Compact Subsets of {Sobolev} Spaces on Manifolds},
  year          = {2020},
  eprint        = {2003.06456},
  archivePrefix = {arXiv},
  primaryClass  = {math.FA},
  doi           = {10.48550/arXiv.2003.06456}
}

@article{shiyao2024rychkov,
  author  = {Shi, Ziming and Yao, Liding},
  title   = {New Estimates of {Rychkov}'s Universal Extension Operator
             for {Lipschitz} Domains and Some Applications},
  journal = {Math. Nachr.},
  volume  = {297},
  number  = {4},
  pages   = {1407--1443},
  year    = {2024},
  doi     = {10.1002/mana.202300047}
}

@misc{berestyckirossi2010principal,
  author        = {Berestycki, Henri and Rossi, Luca},
  title         = {Generalizations and Properties of the Principal
                   Eigenvalue of Elliptic Operators in Unbounded Domains},
  year          = {2010},
  eprint        = {1008.4871},
  archivePrefix = {arXiv},
  primaryClass  = {math.AP},
  doi           = {10.48550/arXiv.1008.4871}
}
```

Skrzypczak and Tintarev study compact Sobolev embeddings on noncompact manifolds using orbital discretizations and quasisymmetric functions. They establish a spotlight lemma that detects compactness through the local vanishing of Sobolev sequences. Their results cover functions supported on domains that become thin at infinity, quasisymmetric functions, and symmetry restrictions. For compact connected groups of isometries, compactness is characterized by the coercivity of the group action.


